% =====================================================================
\chapter{Soluciones de los ejercicios}
\label{ap:soluciones}
% =====================================================================

Cada solución se presenta en cinco partes: datos, planteamiento, desarrollo, resultado y comprobación. En las demostraciones, los datos son las hipótesis, el planteamiento es la estrategia y la comprobación es un caso particular o una verificación independiente. Los valores aproximados se dan con cuatro cifras significativas salvo que se indique otra cosa.

% ---------------------------------------------------------------------
\section{Tema 1}

\begin{solucion}{eje:1-vf}
\datos cinco afirmaciones sobre vectores de $\R^3$.
\planteamiento buscar un contraejemplo o citar el resultado que la demuestra.
\pasos
\begin{enumerate}[label=(\alph*)]
  \item Falsa: $\vu\cdot\vv=\sum u_iv_i$ es un número (\cref{def:escalar}).
  \item Falsa: $\vu=\vi\neq\vcero$ y $\vv=\vj\neq\vcero$ cumplen $\vu\cdot\vv=0$. La condición solo dice que son ortogonales.
  \item Falsa: con $\vu=\vi$ y $\vv=-\vi$, $\norm{\vu+\vv}=0$ y $\norm{\vu}+\norm{\vv}=2$.
  \item Falsa: $\vi\times\vj=\vk$ y $\vj\times\vi=-\vk$ (\cref{prop:vectorial}(i)).
  \item Verdadera: es la \cref{prop:vectorial}(iii).
\end{enumerate}
\resultado \textbf{(a) F; (b) F; (c) F; (d) F; (e) V.}
\comprobacion en (c), la igualdad en la desigualdad triangular exige igualdad en Cauchy--Schwarz con $\vu\cdot\vv\ge0$, es decir, vectores de igual dirección y sentido; $\vi$ y $-\vi$ tienen sentidos opuestos.
\end{solucion}

\begin{solucion}{eje:1-cosenos}
\datos $\va\neq\vcero$ con ángulos $\alpha,\beta,\gamma$ respecto de $\vi,\vj,\vk$; $\va=-\vi+2\vj+2\vk$ en (b).
\planteamiento aplicar la \cref{def:angulo} con los vectores de la base, que son unitarios.
\pasos
\begin{enumerate}[label=\arabic*.]
  \item $\cos\alpha=\dfrac{\va\cdot\vi}{\norm{\va}\norm{\vi}}=\dfrac{a_1}{\norm\va}$, y análogamente $\cos\beta=a_2/\norm\va$, $\cos\gamma=a_3/\norm\va$. Luego $\va/\norm\va=(\cos\alpha,\cos\beta,\cos\gamma)$.
  \item Ese vector es unitario: $\cos^2\alpha+\cos^2\beta+\cos^2\gamma=\norm{\va/\norm\va}^2=1$.
  \item En (b), $\norm\va=\sqrt{1+4+4}=3$, así que $(\cos\alpha,\cos\beta,\cos\gamma)=(-\tfrac13,\tfrac23,\tfrac23)$.
  \item $\alpha=\arccos(-\tfrac13)\approx109{,}47^\circ$ y $\beta=\gamma=\arccos\tfrac23\approx48{,}19^\circ$.
\end{enumerate}
\resultado \textbf{$\va/\norm\va=(\cos\alpha,\cos\beta,\cos\gamma)$; en (b), cosenos $(-\frac13,\frac23,\frac23)$ y ángulos $109{,}47^\circ$, $48{,}19^\circ$, $48{,}19^\circ$.}
\comprobacion $\frac19+\frac49+\frac49=1$. El ángulo con $\vi$ es obtuso porque la primera componente es negativa.
\end{solucion}

\begin{solucion}{eje:1-operaciones}
\datos $\vu=(2,-1,3)$, $\vv=(1,4,-2)$.
\planteamiento operar componente a componente (\cref{def:vector,def:escalar}).
\pasos
\begin{enumerate}[label=\arabic*.]
  \item $\vu+\vv=(3,3,1)$; $3\vu-2\vv=(6-2,\,-3-8,\,9+4)=(4,-11,13)$.
  \item $\vu\cdot\vv=2-4-6=-8$.
  \item $\norm\vu=\sqrt{4+1+9}=\sqrt{14}\approx3{,}742$; $\norm\vv=\sqrt{1+16+4}=\sqrt{21}\approx4{,}583$.
  \item $\vu-\vv=(1,-5,5)$, luego la distancia es $\sqrt{1+25+25}=\sqrt{51}\approx7{,}141$.
  \item Como $\vu\cdot\vv<0$, el ángulo es obtuso.
\end{enumerate}
\resultado \textbf{$(3,3,1)$; $(4,-11,13)$; $-8$; $\sqrt{14}$; $\sqrt{21}$; $\sqrt{51}$; ángulo obtuso.}
\comprobacion por~\eqref{eq:t1-desarrollo}, $\norm{\vu-\vv}^2=14+21-2(-8)=51$.
\end{solucion}

\begin{solucion}{eje:1-unitarios}
\datos tres vectores; en (b), norma 6, dirección $(1,2,-2)/3$ y origen $(-2,0,1)$.
\planteamiento dividir cada vector por su norma; en (b), extremo $=$ origen $+$ vector.
\pasos
\begin{enumerate}[label=\arabic*.]
  \item $\norm{(1,1,1)}=\sqrt3$, unitario $\frac{1}{\sqrt3}(1,1,1)$.
  \item $\norm{(2,-1,2)}=\sqrt{4+1+4}=3$, unitario $\frac13(2,-1,2)$.
  \item $\norm{(3,-6,-2)}=\sqrt{9+36+4}=7$, unitario $\frac17(3,-6,-2)$.
  \item En (b), el vector es $6\cdot\frac13(1,2,-2)=(2,4,-4)$ y el extremo es $(-2,0,1)+(2,4,-4)=(0,4,-3)$.
\end{enumerate}
\resultado \textbf{normas $\sqrt3$, $3$, $7$; unitarios $\frac{(1,1,1)}{\sqrt3}$, $\frac{(2,-1,2)}{3}$, $\frac{(3,-6,-2)}{7}$; extremo $(0,4,-3)$.}
\comprobacion $\norm{(2,4,-4)}=\sqrt{4+16+16}=6$, como se pedía.
\end{solucion}

\begin{solucion}{eje:1-angulos}
\datos (a) $(1,0,-1)$ y $(4,4,-2)$; (b) $(1,1,2)$ y $(2,-1,1)$.
\planteamiento $\cos\theta=\vu\cdot\vv/(\norm\vu\norm\vv)$.
\pasos
\begin{enumerate}[label=\arabic*.]
  \item (a) $\vu\cdot\vv=4+0+2=6$, $\norm\vu=\sqrt2$, $\norm\vv=\sqrt{36}=6$; $\cos\theta=\frac{6}{6\sqrt2}=\frac{1}{\sqrt2}$, $\theta=\pi/4$.
  \item (b) $\vu\cdot\vv=2-1+2=3$, $\norm\vu=\norm\vv=\sqrt6$; $\cos\theta=\frac36=\frac12$, $\theta=\pi/3$.
\end{enumerate}
\resultado \textbf{(a) $45^\circ$; (b) $60^\circ$.}
\comprobacion en (a), $\vu_\perp=\vu-\proj_{\vv}\vu$ tiene norma $\norm\vu\sin45^\circ=1$: $\proj_{\vv}\vu=\frac{6}{36}(4,4,-2)$ y $\vu-\proj_{\vv}\vu=\frac13(1,-2,-2)$, de norma~$1$.
\end{solucion}

\begin{solucion}{eje:1-ortogonal-c}
\datos $\vu=(c,2,-1)$, $\vv=(1,-1,2)$.
\planteamiento el ángulo es recto si $\vu\cdot\vv=0$ y agudo si $\vu\cdot\vv>0$ (\cref{prop:coseno}).
\pasos
\begin{enumerate}[label=\arabic*.]
  \item $\vu\cdot\vv=c-2-2=c-4$.
  \item (a) $c-4=0\iff c=4$. (b) $c-4>0\iff c>4$.
  \item Los vectores nunca son paralelos: $(c,2,-1)=t(1,-1,2)$ exigiría $t=-\frac12$ por la tercera componente y $t=-2$ por la segunda.
\end{enumerate}
\resultado \textbf{(a) $c=4$; (b) $c>4$.}
\comprobacion con $c=4$: $(4,2,-1)\cdot(1,-1,2)=4-2-2=0$.
\end{solucion}

\begin{solucion}{eje:1-proyeccion}
\datos $\mathbf F=(2,-2,1)$, $\va=(3,2,-6)$.
\planteamiento aplicar la \cref{def:proyeccion} y la \cref{prop:descomposicion}.
\pasos
\begin{enumerate}[label=\arabic*.]
  \item $\mathbf F\cdot\va=6-4-6=-4$ y $\norm\va=\sqrt{9+4+36}=7$; la componente escalar es $-\frac47$.
  \item $\proj_{\va}\mathbf F=\frac{-4}{49}(3,2,-6)=\bigl(-\frac{12}{49},-\frac{8}{49},\frac{24}{49}\bigr)$.
  \item $\mathbf F_\perp=\mathbf F-\proj_{\va}\mathbf F=\bigl(\frac{98+12}{49},\frac{-98+8}{49},\frac{49-24}{49}\bigr)=\frac{1}{49}(110,-90,25)$.
\end{enumerate}
\resultado \textbf{componente $-\frac47$; $\mathbf F_\parallel=\frac{1}{49}(-12,-8,24)$; $\mathbf F_\perp=\frac{1}{49}(110,-90,25)$.}
\comprobacion $\mathbf F_\perp\cdot\va=\frac{1}{49}(330-180-150)=0$.
\end{solucion}

\begin{solucion}{eje:1-cubo}
\datos cubo $[0,1]^3$, $P=\vcero$; diagonal principal $(1,1,1)$, diagonal de cara $(1,1,0)$, arista $(1,0,0)$.
\planteamiento \cref{def:angulo}.
\pasos
\begin{enumerate}[label=\arabic*.]
  \item (a) $(1,1,1)\cdot(1,1,0)=2$, normas $\sqrt3$ y $\sqrt2$: $\cos\theta=\frac{2}{\sqrt6}=\sqrt{2/3}$, $\theta\approx35{,}26^\circ$.
  \item (b) $(1,1,1)\cdot(1,0,0)=1$: $\cos\theta=\frac1{\sqrt3}$, $\theta\approx54{,}74^\circ$.
\end{enumerate}
\resultado \textbf{(a) $\arccos\sqrt{2/3}\approx35{,}26^\circ$; (b) $\arccos(1/\sqrt3)\approx54{,}74^\circ$.}
\comprobacion los dos ángulos suman $90^\circ$. La diagonal de cara $(1,1,0)$ es la proyección de $(1,1,1)$ sobre el plano $z=0$, de modo que (a) es el ángulo de la diagonal principal con ese plano, complementario del que forma con $\vk$; y este, por simetría, es igual al de (b). En efecto, $\cos^2$ de uno más $\cos^2$ del otro da $\frac23+\frac13=1$.
\end{solucion}

\begin{solucion}{eje:1-vectorial}
\datos $\va=(1,-2,1)$, $\vb=(2,-1,-1)$; $P=(2,0,1)$, $Q=(3,1,0)$, $R=(-1,1,-1)$.
\planteamiento \cref{def:vectorial}; el área del triángulo es la mitad de $\norm{\overrightarrow{PQ}\times\overrightarrow{PR}}$ (\cref{prop:vectorial}(iv)).
\pasos
\begin{enumerate}[label=\arabic*.]
  \item $\va\times\vb=\bigl((-2)(-1)-1\cdot(-1),\ 1\cdot2-1\cdot(-1),\ 1\cdot(-1)-(-2)\cdot2\bigr)=(3,3,3)$.
  \item $(3,3,3)\cdot\va=3-6+3=0$ y $(3,3,3)\cdot\vb=6-3-3=0$.
  \item $\overrightarrow{PQ}=(1,1,-1)$, $\overrightarrow{PR}=(-3,1,-2)$; su producto vectorial es $\bigl(1\cdot(-2)-(-1)\cdot1,\ (-1)(-3)-1\cdot(-2),\ 1\cdot1-1\cdot(-3)\bigr)=(-1,5,4)$.
  \item $\norm{(-1,5,4)}=\sqrt{1+25+16}=\sqrt{42}$; área $=\frac{\sqrt{42}}{2}\approx3{,}240$.
\end{enumerate}
\resultado \textbf{(a) $(3,3,3)$; (b) $\sqrt{42}/2\approx3{,}240$.}
\comprobacion por la identidad de Lagrange, $\norm{\overrightarrow{PQ}}^2\norm{\overrightarrow{PR}}^2-(\overrightarrow{PQ}\cdot\overrightarrow{PR})^2=3\cdot14-(-3+1+2)^2=42$.
\end{solucion}

\begin{solucion}{eje:1-tetraedro}
\datos $P=(1,0,1)$, $Q=(-1,1,2)$, $R=(0,0,2)$, $S=(3,1,-1)$.
\planteamiento volumen $=\frac16\abs{\overrightarrow{PQ}\cdot(\overrightarrow{PR}\times\overrightarrow{PS})}$ (\cref{prop:volumen}).
\pasos
\begin{enumerate}[label=\arabic*.]
  \item $\overrightarrow{PQ}=(-2,1,1)$, $\overrightarrow{PR}=(-1,0,1)$, $\overrightarrow{PS}=(2,1,-2)$.
  \item $\overrightarrow{PR}\times\overrightarrow{PS}=\bigl(0\cdot(-2)-1\cdot1,\ 1\cdot2-(-1)(-2),\ (-1)\cdot1-0\cdot2\bigr)=(-1,0,-1)$.
  \item $\overrightarrow{PQ}\cdot(-1,0,-1)=2+0-1=1$.
  \item Volumen $=\frac16\cdot1=\frac16$.
\end{enumerate}
\resultado \textbf{$\frac16$.}
\comprobacion desarrollando el determinante de filas $\overrightarrow{PQ},\overrightarrow{PR},\overrightarrow{PS}$ por la primera fila: $-2(0-1)-1(2-2)+1(-1-0)=2-0-1=1$.
\end{solucion}

\begin{solucion}{eje:1-polarizacion}
\datos $\vu,\vv\in\R^n$; en (b), no nulos y no opuestos.
\planteamiento (a) usar~\eqref{eq:t1-desarrollo}; (b) comparar los cosenos de los dos ángulos.
\pasos
\begin{enumerate}[label=\arabic*.]
  \item Por~\eqref{eq:t1-desarrollo}, $\norm{\vu+\vv}^2-\norm{\vu-\vv}^2=(\norm\vu^2+2\vu\cdot\vv+\norm\vv^2)-(\norm\vu^2-2\vu\cdot\vv+\norm\vv^2)=4\,\vu\cdot\vv$. Dividiendo por~4 se obtiene la identidad.
  \item $\vw=\hat\vu+\hat\vv\neq\vcero$, porque $\hat\vu=-\hat\vv$ significaría que $\vu$ y $\vv$ tienen sentidos opuestos.
  \item $\vw\cdot\hat\vu=1+\hat\vu\cdot\hat\vv$ y $\vw\cdot\hat\vv=\hat\vu\cdot\hat\vv+1$, iguales.
  \item El coseno del ángulo entre $\vw$ y $\vu$ es $\vw\cdot\hat\vu/\norm\vw$, y el del ángulo entre $\vw$ y $\vv$ es $\vw\cdot\hat\vv/\norm\vw$; son iguales, y como $\arccos$ es inyectiva en $[-1,1]$, los ángulos coinciden.
\end{enumerate}
\resultado \textbf{$\vu\cdot\vv=\frac14\norm{\vu+\vv}^2-\frac14\norm{\vu-\vv}^2$; el bisector unitario es $(\hat\vu+\hat\vv)/\norm{\hat\vu+\hat\vv}$.}
\comprobacion con $\vu=\vi$, $\vv=3\vj$: $\hat\vu+\hat\vv=(1,1,0)$ forma $45^\circ$ con ambos. Y $\frac14(10)-\frac14(10)=0=\vu\cdot\vv$.
\end{solucion}

\begin{solucion}{eje:1-cs-cuadratica}
\datos $\vu,\vv\in\R^n$, $\vv\neq\vcero$ (si $\vv=\vcero$ la desigualdad es trivial).
\planteamiento un polinomio de grado~2 con coeficiente principal positivo que no toma valores negativos tiene discriminante menor o igual que cero.
\pasos
\begin{enumerate}[label=\arabic*.]
  \item (a) $p(t)=\norm{\vu+t\vv}^2\ge0$ por la positividad del producto escalar.
  \item (b) Por~\eqref{eq:t1-desarrollo} con $t\vv$: $p(t)=\norm\vv^2t^2+2(\vu\cdot\vv)t+\norm\vu^2$, con coeficiente principal $\norm\vv^2>0$.
  \item (c) Si el discriminante fuera positivo, $p$ tendría dos raíces reales distintas y sería negativo entre ellas. Luego $4(\vu\cdot\vv)^2-4\norm\vu^2\norm\vv^2\le0$, es decir, $\abs{\vu\cdot\vv}\le\norm\vu\norm\vv$.
  \item (d) La igualdad equivale a discriminante nulo, es decir, a que $p$ tenga una raíz (doble) $t_0$: $\norm{\vu+t_0\vv}=0$, o sea, $\vu=-t_0\vv$.
\end{enumerate}
\resultado \textbf{$\abs{\vu\cdot\vv}\le\norm\vu\norm\vv$, con igualdad si y solo si $\vu$ es múltiplo de $\vv$.}
\comprobacion con $\vu=(1,2)$, $\vv=(2,4)$: $p(t)=20t^2+20t+5=5(2t+1)^2$, que tiene la raíz doble $t_0=-\frac12$, y $\vu=\frac12\vv$.
\end{solucion}

\begin{solucion}{eje:1-rombo}
\datos paralelogramo de lados (como vectores) $\va$ y $\vb$.
\planteamiento las diagonales son $\va+\vb$ y $\va-\vb$; calcular su producto escalar.
\pasos
\begin{enumerate}[label=\arabic*.]
  \item Con un vértice en el origen, los otros son $\va$, $\vb$ y $\va+\vb$. Las diagonales son $\va+\vb$ (del origen al vértice opuesto) y $\va-\vb$ (de $\vb$ a $\va$).
  \item $(\va+\vb)\cdot(\va-\vb)=\norm\va^2-\va\cdot\vb+\vb\cdot\va-\norm\vb^2=\norm\va^2-\norm\vb^2$.
  \item Por tanto, las diagonales son perpendiculares si y solo si $\norm\va=\norm\vb$.
  \item En un paralelogramo los lados opuestos son iguales; $\norm\va=\norm\vb$ equivale a que los cuatro lados sean iguales.
\end{enumerate}
\resultado \textbf{diagonales perpendiculares $\iff\norm\va=\norm\vb\iff$ rombo.}
\comprobacion para el rombo de lados $\va=(2,0)$, $\vb=(1,\sqrt3)$, las diagonales $(3,\sqrt3)$ y $(1,-\sqrt3)$ cumplen $3-3=0$.
\end{solucion}

\begin{solucion}{eje:1-semicirculo}
\datos circunferencia de centro $O$ y radio $\rho$; $A$, $B$ extremos de un diámetro; $C$ otro punto de la circunferencia.
\planteamiento escribir todo respecto de $O$ y calcular $\overrightarrow{CA}\cdot\overrightarrow{CB}$.
\pasos
\begin{enumerate}[label=\arabic*.]
  \item Sean $\va=\overrightarrow{OA}$ y $\mathbf c=\overrightarrow{OC}$; entonces $\overrightarrow{OB}=-\va$ y $\norm\va=\norm{\mathbf c}=\rho$.
  \item $\overrightarrow{CA}=\va-\mathbf c$ y $\overrightarrow{CB}=-\va-\mathbf c$; ambos son no nulos porque $C\neq A,B$.
  \item $(\va-\mathbf c)\cdot(-\va-\mathbf c)=-\norm\va^2-\va\cdot\mathbf c+\mathbf c\cdot\va+\norm{\mathbf c}^2=\rho^2-\rho^2=0$.
\end{enumerate}
\resultado \textbf{$\overrightarrow{CA}\cdot\overrightarrow{CB}=0$: el ángulo en $C$ es recto.}
\comprobacion con la circunferencia unidad, $A=(1,0)$, $B=(-1,0)$, $C=(0,1)$: $(1,-1)\cdot(-1,-1)=-1+1=0$.
\end{solucion}

\begin{solucion}{eje:1-lagrange-id}
\datos $\vu,\vv\in\R^3$.
\planteamiento desarrollar ambos miembros y agrupar por pares de índices $i<j$.
\pasos
\begin{enumerate}[label=\arabic*.]
  \item Miembro izquierdo: $(u_2v_3-u_3v_2)^2+(u_3v_1-u_1v_3)^2+(u_1v_2-u_2v_1)^2=\sum_{i<j}(u_iv_j-u_jv_i)^2$, una suma sobre los tres pares $(1,2),(1,3),(2,3)$.
  \item Miembro derecho: $\bigl(\sum_iu_i^2\bigr)\bigl(\sum_jv_j^2\bigr)-\bigl(\sum_iu_iv_i\bigr)^2=\sum_{i,j}u_i^2v_j^2-\sum_iu_i^2v_i^2-2\sum_{i<j}u_iv_iu_jv_j$.
  \item Los términos con $i=j$ se cancelan y queda $\sum_{i\neq j}u_i^2v_j^2-2\sum_{i<j}u_iv_iu_jv_j=\sum_{i<j}\bigl(u_i^2v_j^2+u_j^2v_i^2-2u_iv_ju_jv_i\bigr)=\sum_{i<j}(u_iv_j-u_jv_i)^2$.
  \item (b) $\vu\times\vv=\vcero\iff\norm{\vu\times\vv}^2=0\iff(\vu\cdot\vv)^2=\norm\vu^2\norm\vv^2$, que es el caso de igualdad de Cauchy--Schwarz (\cref{teo:cs}): uno es múltiplo del otro.
\end{enumerate}
\resultado \textbf{$\norm{\vu\times\vv}^2=\norm\vu^2\norm\vv^2-(\vu\cdot\vv)^2$; $\vu\times\vv=\vcero$ si y solo si $\vu$ y $\vv$ son paralelos.}
\comprobacion con $\va=(1,-2,1)$, $\vb=(2,-1,-1)$ del \cref{eje:1-vectorial}: $\norm{(3,3,3)}^2=27$ y $6\cdot6-3^2=27$.
\end{solucion}

% ---------------------------------------------------------------------
\section{Tema 2}

\begin{solucion}{eje:2-posicion}
\datos tres pares de planos.
\planteamiento aplicar la \cref{prop:dos-planos}: comparar normales y, si son proporcionales, los términos independientes.
\pasos
\begin{enumerate}[label=(\alph*)]
  \item $(2,4,-2)=2(1,2,-1)$, pero $5\neq2\cdot1$: planos paralelos distintos; sistema incompatible.
  \item $(3,6,-3)=3(1,2,-1)$ y $3=3\cdot1$: el mismo plano; sistema compatible indeterminado (las soluciones forman un plano).
  \item $(1,1,0)$ y $(0,1,1)$ no son proporcionales: se cortan en una recta de dirección $(1,1,0)\times(0,1,1)=(1,-1,1)$; sistema compatible indeterminado.
\end{enumerate}
\resultado \textbf{(a) paralelos, incompatible; (b) coincidentes, indeterminado; (c) se cortan en una recta, indeterminado.}
\comprobacion en (c), el punto $(1,0,1)$ está en ambos planos y $(1,0,1)+t(1,-1,1)$ cumple $x+y=1$ e $y+z=1$ para todo $t$.
\end{solucion}

\begin{solucion}{eje:2-planos}
\datos (a) $P_0=(2,0,-1)$, $\vn=(1,2,-2)$; (b) $P=(1,0,1)$, $Q=(2,-1,2)$, $R=(-1,3,2)$.
\planteamiento \cref{def:plano}; en (b), $\vn=\overrightarrow{PQ}\times\overrightarrow{PR}$.
\pasos
\begin{enumerate}[label=\arabic*.]
  \item (a) $d=\vn\cdot P_0=2+0+2=4$: el plano es $x+2y-2z=4$.
  \item (b) $\overrightarrow{PQ}=(1,-1,1)$, $\overrightarrow{PR}=(-2,3,1)$; $\vn=\bigl((-1)\cdot1-1\cdot3,\ 1\cdot(-2)-1\cdot1,\ 1\cdot3-(-1)(-2)\bigr)=(-4,-3,1)$.
  \item $d=\vn\cdot P=-4+0+1=-3$: $-4x-3y+z=-3$, es decir, $4x+3y-z=3$.
\end{enumerate}
\resultado \textbf{(a) $x+2y-2z=4$; (b) $4x+3y-z=3$.}
\comprobacion en (b): $Q$ da $8-3-2=3$ y $R$ da $-4+9-2=3$.
\end{solucion}

\begin{solucion}{eje:2-diedro}
\datos normales $\vn_1=(2,-1,1)$ y $\vn_2=(1,1,2)$.
\planteamiento el ángulo entre planos es el de sus normales (\cref{prop:dos-planos}).
\pasos
\begin{enumerate}[label=\arabic*.]
  \item $\vn_1\cdot\vn_2=2-1+2=3$; $\norm{\vn_1}=\norm{\vn_2}=\sqrt6$.
  \item $\cos\varphi=\frac36=\frac12$, luego $\varphi=\frac\pi3$.
\end{enumerate}
\resultado \textbf{$60^\circ$.}
\comprobacion la recta de intersección tiene dirección $\vn_1\times\vn_2=(-3,-3,3)$, ortogonal a ambas normales: $-6+3+3=0$ y $-3-3+6=0$.
\end{solucion}

\begin{solucion}{eje:2-interseccion}
\datos (a) punto $(2,-1,-1)$, plano $x-y+2z=3$; (b) recta por $(0,1,2)$ y $(2,0,3)$, plano $x+4y+z=4$.
\planteamiento (a) la normal del plano es el vector director; (b) sustituir la paramétrica en el plano.
\pasos
\begin{enumerate}[label=\arabic*.]
  \item (a) $\vr(t)=(2,-1,-1)+t(1,-1,2)=(2+t,\,-1-t,\,-1+2t)$.
  \item (b) Vector director $(2,0,3)-(0,1,2)=(2,-1,1)$; recta $(2t,\,1-t,\,2+t)$.
  \item Sustituyendo: $2t+4(1-t)+(2+t)=4$, es decir, $6-t=4$, $t=2$.
  \item Punto: $(4,-1,4)$.
\end{enumerate}
\resultado \textbf{(a) $(2+t,-1-t,-1+2t)$; (b) $(4,-1,4)$.}
\comprobacion $4+4(-1)+4=4$; además $(4,-1,4)=(0,1,2)+2(2,-1,1)$ está en la recta.
\end{solucion}

\begin{solucion}{eje:2-gauss}
\datos $x_1-x_2+x_3=2$, $x_2+x_3=0$, $-x_1-x_2+2x_3=3$.
\planteamiento escalonar la matriz ampliada con operaciones (E3) y sustituir hacia atrás; el determinante es el producto de la diagonal final.
\pasos
\begin{enumerate}[label=\arabic*.]
  \item $F_3+F_1$: la tercera fila pasa a $(0,-2,3\mid5)$.
  \item $F_3+2F_2$: $(0,0,5\mid5)$.
  \item $x_3=1$; $x_2=-x_3=-1$; $x_1=2+x_2-x_3=0$.
  \item Solo se han usado operaciones (E3), luego $\det A=1\cdot1\cdot5=5$.
\end{enumerate}
\resultado \textbf{$(x_1,x_2,x_3)=(0,-1,1)$; $\det A=5$.}
\comprobacion $0+1+1=2$, $-1+1=0$, $0+1+2=3$. Desarrollo directo: $1(2+1)+1(0+1)+1(0+1)=5$.
\end{solucion}

\begin{solucion}{eje:2-gauss-casos}
\datos tres sistemas $3\times3$.
\planteamiento eliminación gaussiana (\cref{sec:t2-gauss}).
\pasos
\begin{enumerate}[label=\arabic*.]
  \item (a) $F_2-2F_1\to(0,-3,-1\mid-9)$; $F_3-F_1\to(0,1,-2\mid-4)$; $F_3+\frac13F_2\to(0,0,-\frac73\mid-7)$. Luego $z=3$, $-3y-3=-9\Rightarrow y=2$, $x=6-2-3=1$.
  \item (b) $F_2-F_1\to(0,1,2\mid1)$; $F_3-2F_1\to(0,1,2\mid1)$; $F_3-F_2\to(0,0,0\mid0)$. Con $z$ libre: $y=1-2z$, $x=1-y-z=z$.
  \item (c) Las mismas operaciones llevan la última fila a $(0,0,0\mid1)$: $0=1$, imposible.
\end{enumerate}
\resultado \textbf{(a) compatible determinado, $(1,2,3)$: tres planos con un punto común. (b) compatible indeterminado, $(z,1-2z,z)$: los tres planos comparten una recta. (c) incompatible: los planos se cortan dos a dos en rectas paralelas, sin punto común a los tres.}
\comprobacion (a) $1+2+3=6$, $2-2+3=3$, $1+4-3=2$. (b) con $z=1$, $(1,-1,1)$: $1$, $1-2+3=2$, $2-3+4=3$.
\end{solucion}

\begin{solucion}{eje:2-parametro}
\datos sistema homogéneo con matriz de filas $(1,-1,1)$, $(2,1,1)$, $(-1,c,2)$.
\planteamiento un sistema $3\times3$ homogéneo tiene soluciones no triviales si y solo si $\det A=0$.
\pasos
\begin{enumerate}[label=\arabic*.]
  \item $\det A=1(2-c)-(-1)(4+1)+1(2c+1)=2-c+5+2c+1=c+8$; se anula para $c=-8$.
  \item $\vw=(1,-1,1)\times(2,1,1)=\bigl(-1-1,\ 2-1,\ 1+2\bigr)=(-2,1,3)$ es ortogonal a las dos primeras filas, luego cumple las dos primeras ecuaciones.
  \item Con $c=-8$ las tres filas son coplanarias (\cref{prop:volumen}): la tercera es combinación de las dos primeras, $(-1,-8,2)=5(1,-1,1)-3(2,1,1)$, y por tanto también es ortogonal a $\vw$.
\end{enumerate}
\resultado \textbf{$c=-8$; solución no trivial $(-2,1,3)$ (y sus múltiplos).}
\comprobacion $(-1)(-2)+(-8)(1)+2\cdot3=2-8+6=0$.
\end{solucion}

\begin{solucion}{eje:2-areas}
\datos los triángulos de vértices indicados.
\planteamiento área $=\frac12\abs{\det(\overrightarrow{P_1P_2},\overrightarrow{P_1P_3})}$ (\cref{teo:det-area}).
\pasos
\begin{enumerate}[label=\arabic*.]
  \item (a) Vectores $(1,2)$ y $(1,-1)$: $\det=1\cdot(-1)-2\cdot1=-3$; área $\frac32$.
  \item (b) Desde $(1,2)$: $(0,-3)$ y $(1,1)$; $\det=0\cdot1-(-3)\cdot1=3$; área $\frac32$.
\end{enumerate}
\resultado \textbf{(a) $\frac32$; (b) $\frac32$.}
\comprobacion (a) con base el lado vertical de $(1,2)$ a $(1,-1)$, de longitud $3$, y altura la distancia de $(0,0)$ a la recta $x=1$, que es $1$: $\frac12\cdot3\cdot1=\frac32$.
\end{solucion}

\begin{solucion}{eje:2-inversa}
\datos $A=\begin{pmatrix}4&3\\3&2\end{pmatrix}$, $\vb=(2,3)$.
\planteamiento fórmula~\eqref{eq:t2-inversa} y $\vx=A^{-1}\vb$.
\pasos
\begin{enumerate}[label=\arabic*.]
  \item $\det A=8-9=-1$; $A^{-1}=-1\cdot\begin{pmatrix}2&-3\\-3&4\end{pmatrix}=\begin{pmatrix}-2&3\\3&-4\end{pmatrix}$.
  \item $AA^{-1}=\begin{pmatrix}-8+9&12-12\\-6+6&9-8\end{pmatrix}=I$.
  \item $\vx=A^{-1}\vb=(-4+9,\ 6-12)=(5,-6)$.
\end{enumerate}
\resultado \textbf{$A^{-1}=\begin{pmatrix}-2&3\\3&-4\end{pmatrix}$, $\vx=(5,-6)$.}
\comprobacion $4\cdot5+3(-6)=2$ y $3\cdot5+2(-6)=3$.
\end{solucion}

\begin{solucion}{eje:2-produccion}
\datos columnas de consumo $(1,2,3)$, $(1,3,5)$, $(3,5,8)$; $\vb=(137,279,448)$.
\planteamiento la fila $i$ del sistema cuenta el consumo de $M_i$: $A=\begin{pmatrix}1&1&3\\2&3&5\\3&5&8\end{pmatrix}$, $A\vx=\vb$.
\pasos
\begin{enumerate}[label=\arabic*.]
  \item $\det A=1(24-25)-1(16-15)+3(10-9)=-1-1+3=1\neq0$: solución única.
  \item Eliminación: $F_2-2F_1\to(0,1,-1\mid5)$, $F_3-3F_1\to(0,2,-1\mid37)$, $F_3-2F_2\to(0,0,1\mid27)$.
  \item $x_3=27$, $x_2=5+27=32$, $x_1=137-32-3\cdot27=24$.
  \item (c) Con la columna $(2,5,8)=(1,2,3)+(1,3,5)$, $\det A=1(24-25)-1(16-15)+2(10-9)=0$.
  \item Entonces $A\vx$ siempre está en el plano generado por $(1,2,3)$ y $(1,3,5)$, de normal $(1,2,3)\times(1,3,5)=(1,-2,1)$. Como $(1,-2,1)\cdot\vb=137-558+448=27\neq0$, el sistema es incompatible.
\end{enumerate}
\resultado \textbf{(a) $\det A=1$; (b) $24$, $32$ y $27$ unidades de $P_1$, $P_2$, $P_3$; (c) $\det A=0$ y los consumos dados serían imposibles.}
\comprobacion $24+32+81=137$, $48+96+135=279$, $72+160+216=448$.
\end{solucion}

\begin{solucion}{eje:2-cramer}
\datos $A=\begin{pmatrix}3&-1&1\\-1&2&1\\1&-1&1\end{pmatrix}$, $\vb=(1,2,-3)$.
\planteamiento la incógnita $j$ es $\det A_j/\det A$, donde $A_j$ es $A$ con la columna $j$ sustituida por $\vb$.
\pasos
\begin{enumerate}[label=\arabic*.]
  \item $\det A=3(2+1)-(-1)(-1-1)+1(1-2)=9-2-1=6$.
  \item $\det A_1=\begin{vmatrix}1&-1&1\\2&2&1\\-3&-1&1\end{vmatrix}=1(2+1)+1(2+3)+1(-2+6)=12$, luego $x=2$.
  \item $\det A_2=\begin{vmatrix}3&1&1\\-1&2&1\\1&-3&1\end{vmatrix}=3(2+3)-1(-1-1)+1(3-2)=18$, luego $y=3$.
  \item $\det A_3=\begin{vmatrix}3&-1&1\\-1&2&2\\1&-1&-3\end{vmatrix}=3(-6+2)+1(3-2)+1(1-2)=-12$, luego $z=-2$.
\end{enumerate}
\resultado \textbf{$(x,y,z)=(2,3,-2)$.}
\comprobacion $6-3-2=1$, $-2+6-2=2$, $2-3-2=-3$.
\end{solucion}

\begin{solucion}{eje:2-distancia}
\datos $P_1=(x_1,y_1,z_1)$; plano $\vn\cdot X=d$ con $\vn=(a,b,c)\neq\vcero$.
\planteamiento descomponer $\overrightarrow{P_0P_1}$ en sus partes normal y paralela al plano (\cref{prop:descomposicion}).
\pasos
\begin{enumerate}[label=\arabic*.]
  \item Sea $P_0$ un punto del plano ($\vn\cdot P_0=d$). Para cualquier $X$ del plano, $X-P_0\perp\vn$.
  \item Sea $\mathbf p=\overrightarrow{P_0P_1}=\mathbf p_\parallel+\mathbf p_\perp$ con $\mathbf p_\parallel=\proj_{\vn}\mathbf p$. Entonces $P_1-X=\mathbf p_\parallel+(\mathbf p_\perp-(X-P_0))$, suma de un múltiplo de $\vn$ y de un vector ortogonal a $\vn$.
  \item Por Pitágoras, $\norm{P_1-X}^2=\norm{\mathbf p_\parallel}^2+\norm{\mathbf p_\perp-(X-P_0)}^2\ge\norm{\mathbf p_\parallel}^2$, con igualdad para $X=P_0+\mathbf p_\perp$, que está en el plano.
  \item La distancia es $\norm{\mathbf p_\parallel}=\dfrac{\abs{\vn\cdot(P_1-P_0)}}{\norm\vn}=\dfrac{\abs{ax_1+by_1+cz_1-d}}{\sqrt{a^2+b^2+c^2}}$.
  \item (b) Con $P_1=\vcero$: $\abs d/\norm\vn$. Para $(1,2,3)$ y $2x-y+2z=10$: $\abs{2-2+6-10}/\sqrt{9}=\frac43$.
\end{enumerate}
\resultado \textbf{$d(P_1,\Pi)=\frac{\abs{ax_1+by_1+cz_1-d}}{\sqrt{a^2+b^2+c^2}}$; del origen, $\frac{\abs d}{\norm\vn}$; en el ejemplo, $\frac43$.}
\comprobacion el pie de la perpendicular es $(1,2,3)+\frac49(2,-1,2)=(\frac{17}9,\frac{14}9,\frac{35}9)$, que cumple $\frac{34-14+70}{9}=10$, y está a distancia $\frac49\cdot3=\frac43$.
\end{solucion}

\begin{solucion}{eje:2-vandermonde}
\datos $V=\begin{pmatrix}1&1&1\\x_1&x_2&x_3\\x_1^2&x_2^2&x_3^2\end{pmatrix}$.
\planteamiento restar la primera columna a las otras dos (operación que no cambia el determinante, por la \cref{prop:det-propiedades}(iii) y (iv)) y desarrollar por la primera fila.
\pasos
\begin{enumerate}[label=\arabic*.]
  \item Tras $C_2-C_1$ y $C_3-C_1$, la primera fila es $(1,0,0)$, y $\det V=\begin{vmatrix}x_2-x_1&x_3-x_1\\x_2^2-x_1^2&x_3^2-x_1^2\end{vmatrix}$.
  \item Sacando los factores: $\det V=(x_2-x_1)(x_3-x_1)\begin{vmatrix}1&1\\x_2+x_1&x_3+x_1\end{vmatrix}=(x_2-x_1)(x_3-x_1)(x_3-x_2)$.
  \item Como $(x_1-x_2)(x_2-x_3)=(x_2-x_1)(x_3-x_2)$ (dos cambios de signo), $\det V=(x_1-x_2)(x_2-x_3)(x_3-x_1)$.
  \item La parábola debe cumplir $c_0+c_1x_i+c_2x_i^2=y_i$, $i=1,2,3$: un sistema cuya matriz es $V^{\mathsf T}$, con $\det V^{\mathsf T}=\det V\neq0$ si los $x_i$ son distintos. Por la regla de Cramer en orden~3, la solución existe y es única.
\end{enumerate}
\resultado \textbf{$\det V=(x_1-x_2)(x_2-x_3)(x_3-x_1)$; la parábola por tres puntos de abscisas distintas es única.}
\comprobacion con $(x_1,x_2,x_3)=(0,1,2)$: fórmula, $(-1)(-1)(2)=2$; directamente, $\begin{vmatrix}1&1&1\\0&1&2\\0&1&4\end{vmatrix}=4-2=2$.
\end{solucion}

\begin{solucion}{eje:2-compatibilidad}
\datos $a_1b_2-a_2b_1=0$, $a_1\neq0$.
\planteamiento eliminar $x$ de la segunda ecuación.
\pasos
\begin{enumerate}[label=\arabic*.]
  \item $F_2-\frac{a_2}{a_1}F_1$ da $\bigl(b_2-\frac{a_2b_1}{a_1}\bigr)y=c_2-\frac{a_2}{a_1}c_1$.
  \item El coeficiente de $y$ es $\frac{a_1b_2-a_2b_1}{a_1}=0$, de modo que la ecuación es $0=c_2-\frac{a_2}{a_1}c_1$.
  \item Si $c_2\neq\frac{a_2}{a_1}c_1$, es una contradicción: incompatible. Si $c_2=\frac{a_2}{a_1}c_1$, es $0=0$ y el sistema equivale (\cref{prop:elementales}) a $a_1x+b_1y=c_1$.
  \item Esa ecuación tiene infinitas soluciones: $y=s$ libre y $x=(c_1-b_1s)/a_1$.
\end{enumerate}
\resultado \textbf{compatible $\iff c_2=\frac{a_2}{a_1}c_1$, y entonces con infinitas soluciones (una recta).}
\comprobacion $x+2y=3$, $2x+4y=6$: $c_2=6=2\cdot3$, compatible ($(1,1)$, $(3,0)$, \dots); con $c_2=7$ no hay solución.
\end{solucion}

\begin{solucion}{eje:2-producto-det}
\datos $A=\begin{pmatrix}a&b\\c&d\end{pmatrix}$, $B=\begin{pmatrix}e&f\\g&h\end{pmatrix}$.
\planteamiento (a) cálculo directo; (b) comprobar la definición de inversa; (c) aplicar (a) a $AA^{-1}=I$.
\pasos
\begin{enumerate}[label=\arabic*.]
  \item $AB=\begin{pmatrix}ae+bg&af+bh\\ce+dg&cf+dh\end{pmatrix}$.
  \item $\det(AB)=(ae+bg)(cf+dh)-(af+bh)(ce+dg)$. Los términos $aecf$, $bgdh$ se cancelan con $afce$, $bhdg$, y quedan $aedh+bgcf-afdg-bhce=ad(eh-fg)-bc(eh-fg)=(ad-bc)(eh-fg)$.
  \item (b) $(AB)(B^{-1}A^{-1})=A(BB^{-1})A^{-1}=AIA^{-1}=AA^{-1}=I$, y análogamente $(B^{-1}A^{-1})(AB)=B^{-1}(A^{-1}A)B=I$, usando la asociatividad del producto de matrices.
  \item (c) $\det A\cdot\det(A^{-1})=\det(AA^{-1})=\det I=1$.
\end{enumerate}
\resultado \textbf{$\det(AB)=\det A\det B$; $(AB)^{-1}=B^{-1}A^{-1}$; $\det(A^{-1})=1/\det A$.}
\comprobacion con la matriz del \cref{eje:2-inversa}, $\det A=-1$ y $\det A^{-1}=(-2)(-4)-9=-1=1/(-1)$.
\end{solucion}

% ---------------------------------------------------------------------
\section{Tema 3}

\begin{solucion}{eje:3-conceptos}
\datos tres afirmaciones.
\planteamiento aplicar las definiciones y los ejemplos del tema.
\pasos
\begin{enumerate}[label=(\alph*)]
  \item Falsa: un punto común a los niveles $c_1\neq c_2$ tendría dos valores de $f$.
  \item Falsa: $g(x,y)=\frac{x^2y}{x^4+y^2}$ tiende a~0 por todas las rectas y vale $\frac12$ sobre $y=x^2$ (\cref{ej:t3-no-existe}).
  \item Verdadera: la gráfica son los puntos $(x,y,f(x,y))$ y las curvas de nivel, puntos $(x,y)$ del dominio.
\end{enumerate}
\resultado \textbf{(a) F; (b) F; (c) V.}
\comprobacion el \cref{teo:trayectorias} solo da una condición necesaria, lo que explica (b).
\end{solucion}

\begin{solucion}{eje:3-dominios}
\datos cuatro fórmulas.
\planteamiento raíz: radicando $\ge0$; logaritmo: argumento $>0$; cociente: denominador $\neq0$.
\pasos
\begin{enumerate}[label=(\alph*)]
  \item $x^2+y^2\le1$: disco unidad cerrado. Imagen $[0,1]$ (máximo $1$ en el origen, $0$ en el borde).
  \item $x>y$: semiplano abierto por debajo de la recta $y=x$. Imagen $\R$: sobre $y=0$, $\ln x$ toma todos los valores reales.
  \item $xy\neq0$: el plano sin los dos ejes. Imagen $\R\setminus\{0\}$: con $x=1$, $y=1/k$ se obtiene cualquier $k\neq0$.
  \item $y\ge x^2$: la región sobre la parábola, incluida esta. Imagen $[0,\infty)$: con $x=0$, $\sqrt y$ recorre $[0,\infty)$.
\end{enumerate}
\resultado \textbf{(a) disco cerrado, $[0,1]$; (b) $\{x>y\}$, $\R$; (c) $\{xy\neq0\}$, $\R\setminus\{0\}$; (d) $\{y\ge x^2\}$, $[0,\infty)$.}
\comprobacion (a) $(0{,}6;0{,}8)$ está en el borde y da $\sqrt{1-1}=0$.
\end{solucion}

\begin{solucion}{eje:3-niveles}
\datos cuatro funciones y sus niveles.
\planteamiento resolver $f(x,y)=c$ para cada $c$.
\pasos
\begin{enumerate}[label=(\alph*)]
  \item $1-x-y=c\iff x+y=1-c$: rectas paralelas $x+y=2$, $x+y=1$, $x+y=0$. La gráfica es un plano.
  \item $x^2+y^2=c$: el origen ($c=0$) y circunferencias de radios $1$ y $2$. Paraboloide de vértice el origen, abierto hacia arriba.
  \item $1-x^2-y^2=c\iff x^2+y^2=1-c$: el origen ($c=1$), la circunferencia unidad ($c=0$) y la de radio~$2$ ($c=-3$). Paraboloide abierto hacia abajo con vértice $(0,0,1)$.
  \item $x^2-y^2=1$: hipérbola con vértices $(\pm1,0)$; $c=0$: rectas $y=\pm x$; $c=-1$: hipérbola $y^2-x^2=1$ con vértices $(0,\pm1)$. Silla de montar.
\end{enumerate}
\resultado \textbf{rectas paralelas (plano); circunferencias (paraboloide); circunferencias (paraboloide invertido); hipérbolas y asíntotas (silla).}
\comprobacion (c) $(2,0)$ da $1-4=-3$ y (d) $(1,0)$ da $1$.
\end{solucion}

\begin{solucion}{eje:3-niveles-2}
\datos $f=x^2+4y^2$, $g=y-x^2$, $h=2x+y$, $F=x^2+y^2+z^2$.
\planteamiento identificar la cónica o cuádrica de $f=c$.
\pasos
\begin{enumerate}[label=(\alph*)]
  \item $c>0$: $\frac{x^2}{c}+\frac{y^2}{c/4}=1$, elipses de semiejes $\sqrt c$ y $\sqrt c/2$; $c=0$: el origen; $c<0$: vacío.
  \item $y=x^2+c$: parábolas trasladadas verticalmente.
  \item $2x+y=c$: rectas paralelas de normal $(2,1)$, el vector de coeficientes (\cref{tema:2}). Dos niveles consecutivos $c$ y $c+1$ distan $1/\sqrt5$.
  \item $c>0$: esferas de radio $\sqrt c$ centradas en el origen; $c=0$: el origen.
\end{enumerate}
\resultado \textbf{elipses; parábolas; rectas de normal $(2,1)$; esferas.}
\comprobacion (a) $(\sqrt c,0)$ y $(0,\sqrt c/2)$ dan $c$; (c) la distancia entre $2x+y=c$ y $2x+y=c+1$ es $1/\norm{(2,1)}$ por el \cref{eje:2-distancia}.
\end{solucion}

\begin{solucion}{eje:3-onda}
\datos $F(x,t)=\cos^2(x-ct)$.
\planteamiento $F$ depende solo de $u=x-ct$.
\pasos
\begin{enumerate}[label=\arabic*.]
  \item (a) $F=k$ con $0<k<1$ equivale a $u\in\{\pm\arccos\sqrt k+n\pi:n\in\Z\}$; cada valor $u_0$ da la recta $x-ct=u_0$, es decir, $t=(x-u_0)/c$. Las curvas de nivel son rectas paralelas de pendiente $1/c$ en el plano $(x,t)$.
  \item (b) $F(x,t)=F(x-ct,0)$: la gráfica en el instante $t$ es la del instante~0 trasladada $ct$ hacia la derecha. Con $c=2$ y $t=-1,0,1$, la misma curva $\cos^2x$ aparece desplazada $-2$, $0$ y $2$ unidades.
  \item El perfil se desplaza sin deformarse con velocidad $c$: es una \emph{onda viajera}.
\end{enumerate}
\resultado \textbf{niveles: rectas $x-ct=\text{cte}$; $F$ describe una onda que viaja hacia la derecha con velocidad $c$.}
\comprobacion los máximos ($F=1$) están en $x-2t=n\pi$; en $t=1$, uno de ellos está en $x=2$, dos unidades a la derecha del máximo $x=0$ de $t=0$.
\end{solucion}

\begin{solucion}{eje:3-interseccion-superficies}
\datos $z=x^2-y^2$, $z=2+(x-y)^2$.
\planteamiento igualar las dos expresiones de $z$.
\pasos
\begin{enumerate}[label=\arabic*.]
  \item $x^2-y^2=2+x^2-2xy+y^2\iff2xy-2y^2=2\iff y(x-y)=1$.
  \item Con $y=t\neq0$: $x-t=1/t$, luego $x=t+\frac1t$.
  \item $z=x^2-y^2=t^2+2+\frac1{t^2}-t^2=2+\frac1{t^2}$.
  \item $\vr(t)=\bigl(t+\frac1t,\ t,\ 2+\frac1{t^2}\bigr)$, $t\neq0$; en $t=1$, $\vr(1)=(2,1,3)$.
\end{enumerate}
\resultado \textbf{proyección $y(x-y)=1$; curva $\vr(t)=(t+1/t,\,t,\,2+1/t^2)$, que pasa por $(2,1,3)$ en $t=1$.}
\comprobacion con la segunda superficie, $2+(x-y)^2=2+\frac1{t^2}$, el mismo valor de $z$.
\end{solucion}

\begin{solucion}{eje:3-limites-directos}
\datos cuatro límites.
\planteamiento sustituir cuando la función es continua (\cref{prop:continuas}); si no, simplificar antes.
\pasos
\begin{enumerate}[label=(\alph*)]
  \item Cociente de polinomios con denominador $1+2=3\neq0$: $\frac{2+3}{3}=\frac53$.
  \item Composición de continuas: $e^0\cos0=1$.
  \item $\frac{(x-y)(x+y)}{x-y}=x+y$ para $x\neq y$, con límite $2$.
  \item $x-y=(\sqrt x-\sqrt y)(\sqrt x+\sqrt y)$, luego el cociente es $\sqrt x+\sqrt y\to4$.
\end{enumerate}
\resultado \textbf{(a) $\frac53$; (b) $1$; (c) $2$; (d) $4$.}
\comprobacion (c) en $(1{,}01;\,0{,}99)$: $\frac{1{,}0201-0{,}9801}{0{,}02}=2$.
\end{solucion}

\begin{solucion}{eje:3-no-existe}
\datos (a) $\frac{x^2-y^2}{x^2+y^2}$; (b) $\frac{xy^2}{x^2+y^4}$.
\planteamiento \cref{cor:no-existe}: dos caminos con límites distintos.
\pasos
\begin{enumerate}[label=\arabic*.]
  \item (a) Sobre $y=mx$: $\frac{(1-m^2)x^2}{(1+m^2)x^2}=\frac{1-m^2}{1+m^2}$. Con $m=0$ vale $1$; con $m=1$, $0$.
  \item (b) Sobre $x=my^2$: $\frac{my^4}{m^2y^4+y^4}=\frac{m}{m^2+1}$. Con $m=0$ ($x=0$) vale $0$; con $m=1$, $\frac12$.
\end{enumerate}
\resultado \textbf{ninguno de los dos límites existe.}
\comprobacion en (b), sobre las rectas $y=kx$ el cociente es $\frac{k^2x}{1+k^4x^2}\to0$: las rectas no bastan para detectar la no existencia.
\end{solucion}

\begin{solucion}{eje:3-polares}
\datos cuatro funciones en el origen.
\planteamiento acotar en polares por una función de $r$ que tienda a~0, independiente de $\theta$ (\cref{prop:acotacion}).
\pasos
\begin{enumerate}[label=(\alph*)]
  \item $\frac{r^3\cos^3\theta}{r^2}=r\cos^3\theta$, de valor absoluto $\le r\to0$.
  \item $\frac{r^4\cos^2\theta\sin^2\theta}{r^2}=r^2\cos^2\theta\sin^2\theta\le r^2\to0$.
  \item $\frac{r^2\cos\theta\sin\theta}{r}=r\cos\theta\sin\theta$, de valor absoluto $\le r\to0$.
  \item $r^2\ln r^2=s\ln s$ con $s=r^2\to0^+$; por L'Hôpital, $\frac{\ln s}{1/s}\sim\frac{1/s}{-1/s^2}=-s\to0$. No depende de $\theta$.
\end{enumerate}
\resultado \textbf{los cuatro límites valen $0$.}
\comprobacion (d) en $r=0{,}01$: $10^{-4}\ln10^{-4}\approx-9{,}2\cdot10^{-4}$, próximo a~0.
\end{solucion}

\begin{solucion}{eje:3-continuidad}
\datos (a) $f=\frac{x^3+y^3}{x^2+y^2}$, $f(0,0)=a$; (b) $g=\frac{x^2-y^2}{x^2+y^2}$, $g(0,0)=b$.
\planteamiento $f$ es continua fuera del origen; en el origen hace falta $a=\lim f$.
\pasos
\begin{enumerate}[label=\arabic*.]
  \item (a) En polares, $f=r(\cos^3\theta+\sin^3\theta)$ y $\abs f\le2r\to0$, así que $\lim f=0$.
  \item Por tanto $f$ es continua en el origen si y solo si $a=0$.
  \item (b) Por el \cref{eje:3-no-existe}(a), $g$ no tiene límite en el origen: ningún $b$ la hace continua.
\end{enumerate}
\resultado \textbf{(a) $a=0$; (b) no existe tal $b$.}
\comprobacion en (b), $g(x,0)=1$ y $g(0,y)=-1$, de modo que $b$ tendría que valer a la vez $1$ y $-1$.
\end{solucion}

\begin{solucion}{eje:3-lse}
\datos $f(x,y)=\ln(e^x+e^y)$.
\planteamiento (a) acotar la suma por su mayor sumando; (b) despejar $y$.
\pasos
\begin{enumerate}[label=\arabic*.]
  \item (a) Sin pérdida de generalidad $x\ge y$, de modo que $\max(x,y)=x$ y $e^y\le e^x$. Entonces $e^x\le e^x+e^y\le2e^x$, y como $\ln$ es creciente, $x\le f\le x+\ln2$.
  \item (b) $f=0\iff e^x+e^y=1\iff e^y=1-e^x$, posible solo si $e^x<1$, es decir, $x<0$. Luego $y=\varphi(x)=\ln(1-e^x)$, $x\in(-\infty,0)$.
  \item Cuando $x\to-\infty$, $\varphi(x)\to\ln1=0$: asíntota horizontal $y=0$. Cuando $x\to0^-$, $\varphi(x)\to-\infty$: asíntota vertical $x=0$. La curva es simétrica respecto de $y=x$.
  \item (c) $f(10,0)=\ln(e^{10}+1)\approx10{,}0000454$; $f(1,1)=\ln(2e)=1+\ln2\approx1{,}693$.
\end{enumerate}
\resultado \textbf{$\max(x,y)\le f\le\max(x,y)+\ln2$; nivel~0: $y=\ln(1-e^x)$, $x<0$, con asíntotas $y=0$ y $x=0$; $f(10,0)\approx10{,}00005$ y $f(1,1)\approx1{,}693$.}
\comprobacion la cota $\ln2$ se alcanza cuando $x=y$, como en $f(1,1)-1=\ln2$; cuando un argumento domina, $f$ es casi el máximo: $f(10,0)-10\approx4{,}5\cdot10^{-5}$.
\end{solucion}

\begin{solucion}{eje:3-epsilon-delta}
\datos (a) $f=2x+3y$ en $(1,1)$, $L=5$; (b) $g=\frac{x^2y}{x^2+y^2}$ en $(0,0)$, $L=0$.
\planteamiento acotar $\abs{f-L}$ por un múltiplo de $\norm{\vx-\va}$ con~\eqref{eq:t3-coordenada}.
\pasos
\begin{enumerate}[label=\arabic*.]
  \item (a) $\abs{2x+3y-5}=\abs{2(x-1)+3(y-1)}\le2\abs{x-1}+3\abs{y-1}\le5\norm{(x,y)-(1,1)}$.
  \item Dado $\varepsilon>0$, con $\delta=\varepsilon/5$: si $\norm{(x,y)-(1,1)}<\delta$, entonces $\abs{2x+3y-5}<5\delta=\varepsilon$.
  \item (b) Como $\frac{x^2}{x^2+y^2}\le1$, $\abs g\le\abs y\le\norm{(x,y)}$.
  \item Dado $\varepsilon>0$, con $\delta=\varepsilon$: si $0<\norm{(x,y)}<\delta$, entonces $\abs g<\varepsilon$.
\end{enumerate}
\resultado \textbf{(a) $\delta=\varepsilon/5$; (b) $\delta=\varepsilon$.}
\comprobacion (a) con $\varepsilon=0{,}1$, $\delta=0{,}02$; el punto $(1{,}014;\,1{,}014)$ está a distancia $0{,}0198$ y da $\abs{5{,}07-5}=0{,}07<0{,}1$.
\end{solucion}

\begin{solucion}{eje:3-composicion}
\datos $g$ continua en $\va$; $h$ continua en $b=g(\va)$.
\planteamiento encadenar las dos definiciones $\varepsilon$--$\delta$.
\pasos
\begin{enumerate}[label=\arabic*.]
  \item Sea $\varepsilon>0$. Por la continuidad de $h$ en $b$, existe $\eta>0$ tal que $\abs{s-b}<\eta\Rightarrow\abs{h(s)-h(b)}<\varepsilon$.
  \item Por la continuidad de $g$ en $\va$, aplicada con $\eta$ en el papel de $\varepsilon$, existe $\delta>0$ tal que $\norm{\vx-\va}<\delta\Rightarrow\abs{g(\vx)-b}<\eta$.
  \item Si $\norm{\vx-\va}<\delta$, entonces $s=g(\vx)$ cumple $\abs{s-b}<\eta$, luego $\abs{h(g(\vx))-h(g(\va))}<\varepsilon$.
  \item (b) $g(x,y)=-(x^2+y^2)$ es un polinomio, continuo en $\R^2$, y $h(s)=e^s$ es continua en $\R$.
\end{enumerate}
\resultado \textbf{$h\circ g$ es continua en $\va$; $e^{-(x^2+y^2)}$ es continua en $\R^2$.}
\comprobacion la composición en orden inverso no tiene sentido aquí: $g\circ h$ requeriría que $h$ tomara valores en $\R^2$.
\end{solucion}

\begin{solucion}{eje:3-cuadratica}
\datos $f=x^2+xy+y^2$.
\planteamiento acotar el término cruzado con $\abs{xy}\le\frac12(x^2+y^2)$, que equivale a $(\abs x-\abs y)^2\ge0$.
\pasos
\begin{enumerate}[label=\arabic*.]
  \item $(\abs x-\abs y)^2\ge0\iff x^2+y^2\ge2\abs{xy}$.
  \item Inferior: $f\ge x^2+y^2-\abs{xy}\ge(x^2+y^2)-\frac12(x^2+y^2)=\frac12(x^2+y^2)$.
  \item Superior: $f\le x^2+y^2+\abs{xy}\le\frac32(x^2+y^2)$.
  \item (b) Si $f(x,y)=c$: $\frac12(x^2+y^2)\le c$ da $x^2+y^2\le2c$, y $c\le\frac32(x^2+y^2)$ da $x^2+y^2\ge\frac{2c}{3}$.
  \item Si $(x,y)\neq(0,0)$, $f\ge\frac12(x^2+y^2)>0$.
\end{enumerate}
\resultado \textbf{$\frac12(x^2+y^2)\le f\le\frac32(x^2+y^2)$; los niveles $f=c>0$ están en la corona $\frac{2c}{3}\le x^2+y^2\le2c$.}
\comprobacion las cotas son óptimas: en $(1,1)$, $f=3=\frac32\cdot2$; en $(1,-1)$, $f=1=\frac12\cdot2$.
\end{solucion}

\begin{solucion}{eje:3-trayectorias}
\datos $f=\frac{x^4y^4}{(x^2+y^4)^3}$.
\planteamiento calcular el límite a lo largo de rectas y de la parábola $x=y^2$.
\pasos
\begin{enumerate}[label=\arabic*.]
  \item Sobre los ejes, $f=0$.
  \item Sobre $y=mx$, $m\neq0$: $f=\frac{m^4x^8}{x^6(1+m^4x^2)^3}=\frac{m^4x^2}{(1+m^4x^2)^3}\to0$.
  \item Sobre $x=y^2$: $f=\frac{y^8\cdot y^4}{(2y^4)^3}=\frac{y^{12}}{8y^{12}}=\frac18$.
  \item Los límites $0$ y $\frac18$ son distintos: por el \cref{cor:no-existe}, no existe el límite.
\end{enumerate}
\resultado \textbf{$f\to0$ por toda recta y $f\equiv\frac18$ sobre $x=y^2$; el límite no existe.}
\comprobacion en el punto $(10^{-2},10^{-1})$ de la parábola: $\frac{10^{-8}\cdot10^{-4}}{(2\cdot10^{-4})^3}=\frac{10^{-12}}{8\cdot10^{-12}}=\frac18$.
\end{solucion}

% ---------------------------------------------------------------------
\section{Tema 4}

\begin{solucion}{eje:4-conceptos}
\datos cuatro afirmaciones.
\planteamiento aplicar las \cref{def:parcial,def:direccional} y los teoremas del tema.
\pasos
\begin{enumerate}[label=(\alph*)]
  \item Falsa: $f_x(a,b)$ fija $y=b$ y hace variar $x$; es la pendiente del corte con el plano $y=b$.
  \item Falsa: $\vw$ no es unitario. Con $\vu=\vw/\sqrt2$, $D_{\vu}f=7/\sqrt2\approx4{,}950$.
  \item Falsa: $\frac{xy}{x^2+y^2}$ (con valor~0 en el origen) tiene derivadas parciales nulas en el origen y no es continua allí (\cref{eje:4-no-diferenciable}).
  \item Verdadera: es el \cref{teo:gradiente-normal}.
\end{enumerate}
\resultado \textbf{(a) F; (b) F; (c) F; (d) V.}
\comprobacion en (b), $7/\sqrt2\le\norm{(3,4)}=5$, como exige el \cref{cor:maximo-crecimiento}; el valor $7$ lo violaría.
\end{solucion}

\begin{solucion}{eje:4-parciales}
\datos seis funciones.
\planteamiento derivar respecto de una variable tratando las demás como constantes.
\pasos
\begin{enumerate}[label=(\alph*)]
  \item $f_x=3x^2y-3y^2$; $f_y=x^3-6xy+4y$.
  \item $f_x=2x/y$; $f_y=-x^2/y^2$.
  \item $f_x=3\cos(3x+2y)$; $f_y=2\cos(3x+2y)$.
  \item $f_x=ye^{xy}$; $f_y=xe^{xy}$.
  \item Producto y cadena: $f_x=\ln(2x+y)+\dfrac{2x}{2x+y}$; $f_y=\dfrac{x}{2x+y}$.
  \item $f_x=2xz$; $f_y=-2z^3$; $f_z=x^2-6yz^2$.
\end{enumerate}
\resultado \textbf{las expresiones anteriores.}
\comprobacion (d) $f_{xy}=e^{xy}+xye^{xy}=f_{yx}$, de acuerdo con el teorema de Schwarz.
\end{solucion}

\begin{solucion}{eje:4-schwarz}
\datos (a) $x^my^n$; (b) $\cos(x^2+y)$; (c) $f(x)g(y)$; (d) $w=e^{ax}\sin(ay)$.
\planteamiento calcular las dos derivadas cruzadas por separado.
\pasos
\begin{enumerate}[label=\arabic*.]
  \item (a) $f_x=mx^{m-1}y^n$, $f_{xy}=mnx^{m-1}y^{n-1}$; $f_y=nx^my^{n-1}$, $f_{yx}=mnx^{m-1}y^{n-1}$.
  \item (b) $f_x=-2x\sin(x^2+y)$, $f_{xy}=-2x\cos(x^2+y)$; $f_y=-\sin(x^2+y)$, $f_{yx}=-2x\cos(x^2+y)$.
  \item (c) $(f(x)g(y))_x=f'(x)g(y)$, y su derivada respecto de $y$ es $f'(x)g'(y)$; en el otro orden, $f(x)g'(y)$ y después $f'(x)g'(y)$.
  \item (d) $w_{xx}=a^2e^{ax}\sin(ay)$ y $w_{yy}=-a^2e^{ax}\sin(ay)$; su suma es~0.
\end{enumerate}
\resultado \textbf{en los tres casos $f_{xy}=f_{yx}$; $w$ es armónica.}
\comprobacion (a) con $m=2$, $n=3$: $x^2y^3$ da $f_{xy}=6xy^2$ por ambos caminos.
\end{solucion}

\begin{solucion}{eje:4-tangente}
\datos (a) $z=xy^2$ en $(1,1,1)$; (b) $z=e^x\cos y$ en $(0,0,1)$.
\planteamiento \cref{def:tangente}.
\pasos
\begin{enumerate}[label=\arabic*.]
  \item (a) $z_x=y^2=1$, $z_y=2xy=2$: $z=1+(x-1)+2(y-1)=x+2y-2$.
  \item (b) $z_x=e^x\cos y=1$, $z_y=-e^x\sin y=0$: $z=1+x$.
  \item Aproximación: $e^{0{,}1}\cos0{,}1\approx1+0{,}1=1{,}1$. Valor exacto: $1{,}10517\cdot0{,}99500\approx1{,}09965$.
\end{enumerate}
\resultado \textbf{(a) $z=x+2y-2$; (b) $z=1+x$; aproximación $1{,}1$ frente a $1{,}0996$.}
\comprobacion el error, $3{,}5\cdot10^{-4}$, es mucho menor que los incrementos ($0{,}1$), como corresponde a un error de orden superior al primero.
\end{solucion}

\begin{solucion}{eje:4-aproximacion}
\datos (a) $c=\sqrt{a^2+b^2}$, $a=3$, $b=4$, $\Delta a=\Delta b=0{,}01$; (b) $V=xyz$, $(x,y,z)=(5,10,20)$, $\abs{\Delta}\le0{,}1$.
\planteamiento $\Delta f\approx f_x\Delta x+f_y\Delta y\,(+f_z\Delta z)$.
\pasos
\begin{enumerate}[label=\arabic*.]
  \item (a) $c=5$, $c_a=a/c=0{,}6$, $c_b=b/c=0{,}8$: $\Delta c\approx0{,}6\cdot0{,}01+0{,}8\cdot0{,}01=0{,}014$.
  \item (b) $V_x=yz=200$, $V_y=xz=100$, $V_z=xy=50$.
  \item El error máximo se da con los tres errores del mismo signo: $\abs{dV}\le(200+100+50)\cdot0{,}1=35$~cm$^3$.
\end{enumerate}
\resultado \textbf{(a) $\Delta c\approx0{,}014$; (b) $\pm35$~cm$^3$ (un 3{,}5\,\% de $1000$~cm$^3$).}
\comprobacion (a) $\sqrt{3{,}01^2+4{,}01^2}-5\approx0{,}0140004$; (b) $5{,}1\cdot10{,}1\cdot20{,}1-1000\approx35{,}35$.
\end{solucion}

\begin{solucion}{eje:4-gradiente}
\datos (a) $f=x^3+2y^3$, $P=(1,1)$, $\va=(1,-1)$; (b) $w=xy/z$, $P=(2,-1,1)$, $\va=(1,2,-2)$.
\planteamiento $D_{\vu}f=\nabla f\cdot\vu$ con $\vu=\va/\norm\va$ (\cref{teo:direccional}).
\pasos
\begin{enumerate}[label=\arabic*.]
  \item (a) $\nabla f=(3x^2,6y^2)=(3,6)$; $\vu=(1,-1)/\sqrt2$; $D_{\vu}f=(3-6)/\sqrt2=-3/\sqrt2\approx-2{,}121$.
  \item (b) $\nabla w=(y/z,\ x/z,\ -xy/z^2)=(-1,2,2)$; $\norm\va=3$, $\vu=(1,2,-2)/3$; $D_{\vu}w=(-1+4-4)/3=-\frac13$.
\end{enumerate}
\resultado \textbf{(a) $\nabla f=(3,6)$, $D_{\vu}f=-\frac{3}{\sqrt2}$; (b) $\nabla w=(-1,2,2)$, $D_{\vu}w=-\frac13$.}
\comprobacion (a) $f(1+t/\sqrt2,1-t/\sqrt2)$ con $t=0{,}001$: $\frac{f-3}{t}\approx-2{,}12$.
\end{solucion}

\begin{solucion}{eje:4-max-direccional}
\datos $w=\ln(4x-3y)$, punto $(1,1)$.
\planteamiento \cref{cor:maximo-crecimiento}.
\pasos
\begin{enumerate}[label=\arabic*.]
  \item $\nabla w=\bigl(\frac{4}{4x-3y},\frac{-3}{4x-3y}\bigr)=(4,-3)$ en $(1,1)$; $\norm{\nabla w}=5$.
  \item Máximo $5$ en la dirección $(4,-3)/5$; mínimo $-5$ en $(-4,3)/5$.
  \item Nula en las direcciones perpendiculares: $\pm(3,4)/5$.
\end{enumerate}
\resultado \textbf{máximo $5$ según $(\frac45,-\frac35)$; mínimo $-5$ según $(-\frac45,\frac35)$; nula según $\pm(\frac35,\frac45)$.}
\comprobacion $(4,-3)\cdot(3,4)=0$; además, $w$ es constante sobre las rectas $4x-3y=\text{cte}$, de dirección $(3,4)$.
\end{solucion}

\begin{solucion}{eje:4-tangente-nivel}
\datos (a) $F=xy^2z^3=12$, $P_0=(3,2,1)$; (b) $G=x^2+4y^2+9z^2=14$, $P_0=(1,1,1)$.
\planteamiento plano de normal $\nabla F(P_0)$ (\cref{teo:gradiente-normal}).
\pasos
\begin{enumerate}[label=\arabic*.]
  \item (a) $\nabla F=(y^2z^3,\ 2xyz^3,\ 3xy^2z^2)=(4,12,36)$. Plano: $4(x-3)+12(y-2)+36(z-1)=0$, es decir, $x+3y+9z=18$.
  \item (b) $\nabla G=(2x,8y,18z)=(2,8,18)$. Plano: $2(x-1)+8(y-1)+18(z-1)=0$, es decir, $x+4y+9z=14$.
\end{enumerate}
\resultado \textbf{(a) $x+3y+9z=18$; (b) $x+4y+9z=14$.}
\comprobacion los puntos de tangencia cumplen las ecuaciones: $3+6+9=18$ y $1+4+9=14$.
\end{solucion}

\begin{solucion}{eje:4-cadena-t}
\datos (a) $w=xyz$, $(t,t^2,t^3)$; (b) $w=x^2-y^2$, $(\cos t,\sin t)$.
\planteamiento \cref{teo:cadena} y comparación con la sustitución directa.
\pasos
\begin{enumerate}[label=\arabic*.]
  \item (a) $w_xx'+w_yy'+w_zz'=yz\cdot1+xz\cdot2t+xy\cdot3t^2=t^5+2t^5+3t^5=6t^5$.
  \item Directamente: $w=t^6$, $w'=6t^5$.
  \item (b) $2x(-\sin t)-2y\cos t=-2\cos t\sin t-2\sin t\cos t=-2\sin2t$.
  \item Directamente: $w=\cos^2t-\sin^2t=\cos2t$, $w'=-2\sin2t$.
\end{enumerate}
\resultado \textbf{(a) $6t^5$; (b) $-2\sin2t$.}
\comprobacion los dos métodos coinciden en ambos casos.
\end{solucion}

\begin{solucion}{eje:4-cadena-uv}
\datos $\nabla f(0,1)=(2,3)$; $x=u^2-v^2$, $y=uv$; $(u,v)=(1,1)$.
\planteamiento fórmula~\eqref{eq:t4-cadena2} en el punto $(x,y)=(0,1)$ que corresponde a $(1,1)$.
\pasos
\begin{enumerate}[label=\arabic*.]
  \item En $(u,v)=(1,1)$: $x=0$, $y=1$, así que se usa $\nabla f(0,1)=(2,3)$.
  \item $x_u=2u=2$, $x_v=-2v=-2$, $y_u=v=1$, $y_v=u=1$.
  \item $w_u=2\cdot2+3\cdot1=7$; $w_v=2\cdot(-2)+3\cdot1=-1$.
\end{enumerate}
\resultado \textbf{$w_u=7$, $w_v=-1$.}
\comprobacion con $f(x,y)=2x+3y$ (que tiene ese gradiente): $w=2u^2-2v^2+3uv$, $w_u=4u+3v=7$ y $w_v=-4v+3u=-1$.
\end{solucion}

\begin{solucion}{eje:4-softmax}
\datos $f=\ln(e^x+e^y)$.
\planteamiento derivar con la regla de la cadena; en (b), \cref{teo:direccional}; en (c), \cref{def:tangente}.
\pasos
\begin{enumerate}[label=\arabic*.]
  \item (a) $f_x=\frac{e^x}{e^x+e^y}=\sigma_1$ y $f_y=\sigma_2$. Son cocientes de números positivos con el numerador menor que el denominador, luego están en $(0,1)$, y $\sigma_1+\sigma_2=1$.
  \item (b) $D_{\vu}f=(\sigma_1+\sigma_2)/\sqrt2=1/\sqrt2$ en todo punto.
  \item En efecto, $f(x+t,y+t)=\ln\bigl(e^t(e^x+e^y)\bigr)=t+f(x,y)$: al avanzar $t$ en cada coordenada (una longitud $t\sqrt2$), $f$ crece $t$.
  \item (c) En $(1,0)$: $f=\ln(e+1)\approx1{,}3133$, $\nabla f\approx(0{,}7311;\,0{,}2689)$. Con $(\Delta x,\Delta y)=(0{,}1;\,-0{,}1)$: $f\approx1{,}3133+0{,}07311-0{,}02689=1{,}3595$.
  \item Valor exacto: $\ln(e^{1{,}1}+e^{-0{,}1})\approx1{,}3633$.
\end{enumerate}
\resultado \textbf{$\nabla f=\softmax(x,y)$; $D_{\vu}f=1/\sqrt2$; aproximación $1{,}3595$ frente a $1{,}3633$.}
\comprobacion el error, $0{,}004$, es del orden de $\norm{\Delta}^2=0{,}02$ por la curvatura de $f$, como corresponde a una aproximación lineal.
\end{solucion}

\begin{solucion}{eje:4-polares}
\datos $w=f(x,y)$, $x=r\cos\theta$, $y=r\sin\theta$.
\planteamiento calcular $w_r$ y $w_\theta$ con~\eqref{eq:t4-cadena2} y sumar sus cuadrados.
\pasos
\begin{enumerate}[label=\arabic*.]
  \item $x_r=\cos\theta$, $y_r=\sin\theta$, $x_\theta=-r\sin\theta$, $y_\theta=r\cos\theta$.
  \item $w_r=w_x\cos\theta+w_y\sin\theta$; $\frac1rw_\theta=-w_x\sin\theta+w_y\cos\theta$.
  \item $w_r^2=w_x^2\cos^2\theta+2w_xw_y\sin\theta\cos\theta+w_y^2\sin^2\theta$.
  \item $\frac{1}{r^2}w_\theta^2=w_x^2\sin^2\theta-2w_xw_y\sin\theta\cos\theta+w_y^2\cos^2\theta$.
  \item Sumando, los términos cruzados se cancelan y $\cos^2\theta+\sin^2\theta=1$: $w_r^2+\frac{1}{r^2}w_\theta^2=w_x^2+w_y^2$.
\end{enumerate}
\resultado \textbf{$w_x^2+w_y^2=w_r^2+r^{-2}w_\theta^2$.}
\comprobacion con $w=x^2+y^2=r^2$: izquierda $4x^2+4y^2=4r^2$; derecha $(2r)^2+0=4r^2$ (\cref{ej:t4-polares}).
\end{solucion}

\begin{solucion}{eje:4-euler}
\datos (a) $w=f(y/x)$; (b) $g(tx,ty)=t^kg(x,y)$; (c) $g=x^2y+y^3$.
\planteamiento (a) regla de la cadena de una variable; (b) derivar la identidad respecto de $t$ y tomar $t=1$.
\pasos
\begin{enumerate}[label=\arabic*.]
  \item (a) $w_x=f'(y/x)\cdot(-y/x^2)$ y $w_y=f'(y/x)\cdot(1/x)$; así $xw_x+yw_y=f'(y/x)\bigl(-\frac yx+\frac yx\bigr)=0$.
  \item (b) Sea $\varphi(t)=g(tx,ty)$. Por la regla de la cadena, $\varphi'(t)=xg_x(tx,ty)+yg_y(tx,ty)$.
  \item Por la homogeneidad, $\varphi(t)=t^kg(x,y)$, luego $\varphi'(t)=kt^{k-1}g(x,y)$.
  \item Igualando en $t=1$: $xg_x+yg_y=kg$.
  \item (c) $g$ es homogénea de grado~3; $xg_x+yg_y=x(2xy)+y(x^2+3y^2)=3x^2y+3y^3=3g$.
\end{enumerate}
\resultado \textbf{$xw_x+yw_y=0$; $xg_x+yg_y=kg$; para $g=x^2y+y^3$, $k=3$.}
\comprobacion (a) es el caso $k=0$ de (b): $f(y/x)$ es homogénea de grado~0.
\end{solucion}

\begin{solucion}{eje:4-material}
\datos $f(x,y,z,t)$ diferenciable; $\vr(t)=(x(t),y(t),z(t))$.
\planteamiento regla de la cadena con cuatro variables intermedias, $x(t)$, $y(t)$, $z(t)$ y $s(t)=t$.
\pasos
\begin{enumerate}[label=\arabic*.]
  \item La función compuesta es $F(t)=f(x(t),y(t),z(t),s(t))$ con $s(t)=t$, $s'(t)=1$.
  \item Por el \cref{teo:cadena} (en cuatro variables), $F'(t)=f_xx'+f_yy'+f_zz'+f_t\cdot1$.
  \item Los tres primeros sumandos son $\nabla f\cdot\vr\,'(t)=\vv\cdot\nabla f$.
\end{enumerate}
\resultado \textbf{$\dfrac{Df}{Dt}=\dfrac{\partial f}{\partial t}+\vv\cdot\nabla f$.}
\comprobacion con $f=x+t$ y $x(t)=t^2$: $F=t^2+t$, $F'=2t+1$; la fórmula da $1+2t\cdot1$.
\end{solucion}

\begin{solucion}{eje:4-no-diferenciable}
\datos $f=\frac{xy}{x^2+y^2}$, $f(0,0)=0$.
\planteamiento (a) definición de derivada parcial; (b) \cref{teo:diferenciable}(i); (c) cálculo directo.
\pasos
\begin{enumerate}[label=\arabic*.]
  \item (a) $f(h,0)=0$ para todo $h$, luego $\frac{f(h,0)-f(0,0)}{h}=0\to0$: $f_x(0,0)=0$. Por simetría, $f_y(0,0)=0$.
  \item (b) Sobre $y=x$, $f=\frac12$, de modo que $f$ no tiende a $f(0,0)=0$: no es continua en el origen y, por el \cref{teo:diferenciable}, no es diferenciable.
  \item (c) $f(t\vu)=\frac{t^2/2}{t^2}=\frac12$ para $t\neq0$; el cociente es $\frac{1}{2t}$, que no tiene límite finito cuando $t\to0$.
  \item Si $f$ fuera diferenciable, el \cref{teo:direccional} daría $D_{\vu}f(\vcero)=\nabla f(\vcero)\cdot\vu=0$; la falta de diferenciabilidad explica la discrepancia.
\end{enumerate}
\resultado \textbf{$f_x(0,0)=f_y(0,0)=0$; $f$ no es diferenciable (ni continua) en el origen; $D_{\vu}f(\vcero)$ no existe.}
\comprobacion la existencia de parciales solo usa los ejes, donde $f\equiv0$; la diagonal, donde $f\equiv\frac12$, revela el problema.
\end{solucion}

% ---------------------------------------------------------------------
\section{Tema 5}

\begin{solucion}{eje:5-conceptos}
\datos cuatro afirmaciones.
\planteamiento contraejemplos.
\pasos
\begin{enumerate}[label=(\alph*)]
  \item Falsa: $x^2-y^2$ en el origen es un punto de silla.
  \item Falsa: $x^4+y^4$ tiene $\Delta=0$ en el origen y un mínimo (\cref{eje:5-degenerado}).
  \item Falsa: las soluciones son candidatos. En el \cref{eje:5-caja}, el sistema tiene también soluciones con volumen $0$, que no son el máximo, y en general un candidato puede no ser extremo.
  \item Falsa: $f(x,y)=x$ en el disco abierto no alcanza su supremo~$1$, porque el punto $(1,0)$ no pertenece al disco. Weierstrass exige un conjunto cerrado.
\end{enumerate}
\resultado \textbf{las cuatro son falsas.}
\comprobacion en (d), $f$ sí alcanza el máximo $1$ en el disco cerrado, en $(1,0)$.
\end{solucion}

\begin{solucion}{eje:5-clasificar-1}
\datos (a) $f=x^2-xy-2y^2-3x-3y+1$; (b) $g=3x^2+xy+y^2-x-2y+4$.
\planteamiento resolver $\nabla f=\vcero$ y aplicar el \cref{teo:hessiano}.
\pasos
\begin{enumerate}[label=\arabic*.]
  \item (a) $f_x=2x-y-3=0$, $f_y=-x-4y-3=0$. De la primera, $y=2x-3$; sustituyendo, $-x-8x+12-3=0$, $x=1$, $y=-1$.
  \item $H_f=\begin{pmatrix}2&-1\\-1&-4\end{pmatrix}$, $\Delta=-8-1=-9<0$: punto de silla.
  \item (b) $g_x=6x+y-1=0$, $g_y=x+2y-2=0$. De la primera, $y=1-6x$; sustituyendo, $x+2-12x-2=0$, $x=0$, $y=1$.
  \item $H_g=\begin{pmatrix}6&1\\1&2\end{pmatrix}$, $\Delta=12-1=11>0$, $g_{xx}=6>0$: mínimo local, $g(0,1)=1-2+4=3$.
\end{enumerate}
\resultado \textbf{(a) silla en $(1,-1)$; (b) mínimo en $(0,1)$ con valor~$3$.}
\comprobacion (b) $g(0{,}1;\,1)=0{,}03+0{,}1+1-0{,}1-2+4=3{,}03>3$ y $g(0;\,1{,}1)=1{,}21-2{,}2+4=3{,}01>3$.
\end{solucion}

\begin{solucion}{eje:5-clasificar-2}
\datos $f=x^3-3xy+y^3$.
\planteamiento resolver $\nabla f=\vcero$ y aplicar el \cref{teo:hessiano}.
\pasos
\begin{enumerate}[label=\arabic*.]
  \item $f_x=3x^2-3y=0\Rightarrow y=x^2$; $f_y=-3x+3y^2=0\Rightarrow x=y^2$.
  \item Sustituyendo, $x=x^4$, es decir, $x(x^3-1)=0$: $x=0$ o $x=1$. Puntos críticos $(0,0)$ y $(1,1)$.
  \item $H_f=\begin{pmatrix}6x&-3\\-3&6y\end{pmatrix}$. En $(0,0)$, $\Delta=-9<0$: silla. En $(1,1)$, $\Delta=36-9=27>0$ y $f_{xx}=6>0$: mínimo local, $f(1,1)=-1$.
\end{enumerate}
\resultado \textbf{silla en $(0,0)$; mínimo local en $(1,1)$, de valor $-1$.}
\comprobacion el mínimo no es absoluto: $f(x,0)=x^3\to-\infty$ cuando $x\to-\infty$.
\end{solucion}

\begin{solucion}{eje:5-clasificar-3}
\datos $f=2x^4+y^2-xy+1$.
\planteamiento resolver $\nabla f=\vcero$ y aplicar el \cref{teo:hessiano}.
\pasos
\begin{enumerate}[label=\arabic*.]
  \item $f_y=2y-x=0\Rightarrow y=x/2$; $f_x=8x^3-y=0\Rightarrow8x^3-x/2=0\Rightarrow x(16x^2-1)=0$.
  \item Puntos críticos: $(0,0)$, $(\frac14,\frac18)$ y $(-\frac14,-\frac18)$.
  \item $H_f=\begin{pmatrix}24x^2&-1\\-1&2\end{pmatrix}$. En $(0,0)$: $\Delta=-1<0$, silla.
  \item En $(\pm\frac14,\pm\frac18)$: $24x^2=\frac32$, $\Delta=3-1=2>0$, $f_{xx}>0$: mínimos locales.
  \item $f(\frac14,\frac18)=2\cdot\frac{1}{256}+\frac1{64}-\frac1{32}+1=\frac{2+4-8}{256}+1=\frac{127}{128}$, y lo mismo en el simétrico.
\end{enumerate}
\resultado \textbf{silla en $(0,0)$; mínimos en $\pm(\frac14,\frac18)$ con valor $\frac{127}{128}$.}
\comprobacion $f(-x,-y)=f(x,y)$, lo que explica la simetría de los mínimos; y $\frac{127}{128}<1=f(0,0)$.
\end{solucion}

\begin{solucion}{eje:5-openstax}
\datos (a) $f=x^3+2xy-2x-4y$; (b) $g=4x^2+9y^2+8x-36y+24$.
\planteamiento \cref{teo:hessiano}; en (b), completar cuadrados.
\pasos
\begin{enumerate}[label=\arabic*.]
  \item (a) $f_y=2x-4=0\Rightarrow x=2$; $f_x=3x^2+2y-2=0\Rightarrow y=(2-12)/2=-5$. Punto $(2,-5)$.
  \item $H_f=\begin{pmatrix}6x&2\\2&0\end{pmatrix}$, $\Delta=-4<0$: silla, con $f(2,-5)=8-20-4+20=4$.
  \item (b) $g_x=8x+8=0$, $g_y=18y-36=0$: punto $(-1,2)$. $H_g=\operatorname{diag}(8,18)$, $\Delta=144>0$, $g_{xx}>0$: mínimo.
  \item $g=4(x+1)^2-4+9(y-2)^2-36+24=4(x+1)^2+9(y-2)^2-16$: mínimo absoluto $-16$ en $(-1,2)$.
\end{enumerate}
\resultado \textbf{(a) silla en $(2,-5)$; (b) mínimo absoluto $-16$ en $(-1,2)$.}
\comprobacion (b) $g(-1,2)=4+36-8-72+24=-16$.
\end{solucion}

\begin{solucion}{eje:5-absolutos}
\datos $f=xy-x-y+2=(x-1)(y-1)+1$ en $[0,2]^2$.
\planteamiento puntos críticos interiores, frontera y comparación (\cref{teo:weierstrass}).
\pasos
\begin{enumerate}[label=\arabic*.]
  \item Interior: $f_x=y-1$, $f_y=x-1$; punto crítico $(1,1)$, $f=1$ (silla: $\Delta=-1$).
  \item Lado $x=0$: $f=2-y\in[0,2]$. Lado $x=2$: $f=y\in[0,2]$. Lado $y=0$: $f=2-x$. Lado $y=2$: $f=x$.
  \item Valores candidatos: $1$ y, en los vértices, $f(0,0)=2$, $f(2,2)=2$, $f(0,2)=0$, $f(2,0)=0$.
\end{enumerate}
\resultado \textbf{máximo $2$ en $(0,0)$ y $(2,2)$; mínimo $0$ en $(0,2)$ y $(2,0)$.}
\comprobacion $(x-1)(y-1)\in[-1,1]$ en el cuadrado, luego $f\in[0,2]$.
\end{solucion}

\begin{solucion}{eje:5-disco}
\datos $f=x+\sqrt3\,y+2$ en $x^2+y^2\le1$.
\planteamiento no hay puntos críticos interiores ($\nabla f=(1,\sqrt3)\neq\vcero$); estudiar la frontera.
\pasos
\begin{enumerate}[label=\arabic*.]
  \item Frontera: $x=\cos t$, $y=\sin t$: $f=\cos t+\sqrt3\sin t+2=2\cos(t-\frac\pi3)+2$.
  \item Máximo $4$ para $t=\frac\pi3$, en $(\frac12,\frac{\sqrt3}2)$; mínimo $0$ para $t=\frac{4\pi}3$, en $(-\frac12,-\frac{\sqrt3}2)$.
\end{enumerate}
\resultado \textbf{máximo $4$ en $(\frac12,\frac{\sqrt3}2)$; mínimo $0$ en $(-\frac12,-\frac{\sqrt3}2)$.}
\comprobacion por Cauchy--Schwarz, $\abs{(1,\sqrt3)\cdot(x,y)}\le2\cdot1$, luego $0\le f\le4$, con igualdad en los múltiplos unitarios de $(1,\sqrt3)$.
\end{solucion}

\begin{solucion}{eje:5-caja}
\datos $V=xyz$, $g=x+2y+3z=18$, $x,y,z\ge0$.
\planteamiento Lagrange en el interior del triángulo $T$ que el plano determina en el primer octante.
\pasos
\begin{enumerate}[label=\arabic*.]
  \item $T$ es cerrado y acotado y $V$ es continua: el máximo existe. En el borde de $T$ alguna coordenada es nula y $V=0$; en el interior, $V>0$. Luego el máximo está en el interior.
  \item Ecuaciones: $yz=\lambda$, $xz=2\lambda$, $xy=3\lambda$.
  \item Multiplicando por $x$, $y$, $z$ respectivamente: $xyz=\lambda x=2\lambda y=3\lambda z$. Como $V>0$, $\lambda\neq0$, y $x=2y=3z$.
  \item Sustituyendo en la restricción: $x+x+x=18$, $x=6$, $y=3$, $z=2$, $V=36$.
\end{enumerate}
\resultado \textbf{$P=(6,3,2)$, volumen máximo $36$.}
\comprobacion otro punto del plano, $(3,3,3)$, da $V=27<36$; y $6+6+6=18$.
\end{solucion}

\begin{solucion}{eje:5-hiperbola}
\datos $f=x^2-y^2$, $g=x^2+y^2=3$.
\planteamiento resolver $\nabla f=\lambda\nabla g$ sin dividir por variables que puedan anularse.
\pasos
\begin{enumerate}[label=\arabic*.]
  \item Ecuaciones: $2x=2\lambda x$ y $-2y=2\lambda y$, es decir, $x(1-\lambda)=0$ e $y(1+\lambda)=0$.
  \item Si $\lambda=1$: $y=0$ y $x=\pm\sqrt3$; $f=3$.
  \item Si $\lambda=-1$: $x=0$ e $y=\pm\sqrt3$; $f=-3$.
  \item Si $\lambda\neq\pm1$: $x=y=0$, que no está en la circunferencia.
  \item La circunferencia es compacta: máximo $3$ y mínimo $-3$.
  \item Con $\lambda=1$, $\nabla f=(2x,0)=\nabla g$: mismo sentido. Con $\lambda=-1$, $\nabla f=(0,-2y)=-\nabla g$: sentidos opuestos.
\end{enumerate}
\resultado \textbf{máximo $3$ en $(\pm\sqrt3,0)$ con $\lambda=1$; mínimo $-3$ en $(0,\pm\sqrt3)$ con $\lambda=-1$.}
\comprobacion en polares, $f=3\cos2\theta\in[-3,3]$.
\end{solucion}

\begin{solucion}{eje:5-resistencias}
\datos $P=I_1^2R_1+I_2^2R_2$, $I_1+I_2=I$.
\planteamiento Lagrange (o sustitución).
\pasos
\begin{enumerate}[label=\arabic*.]
  \item $2I_1R_1=\lambda$ y $2I_2R_2=\lambda$, luego $I_1R_1=I_2R_2$.
  \item Con $I_2=I-I_1$: $I_1R_1=(I-I_1)R_2$, $I_1=\dfrac{IR_2}{R_1+R_2}$ e $I_2=\dfrac{IR_1}{R_1+R_2}$.
  \item Es un mínimo: sustituyendo, $P(I_1)=I_1^2R_1+(I-I_1)^2R_2$ es una parábola con $P''=2(R_1+R_2)>0$.
\end{enumerate}
\resultado \textbf{$I_1=\frac{R_2}{R_1+R_2}I$, $I_2=\frac{R_1}{R_1+R_2}I$, que cumplen $I_1R_1=I_2R_2$ (igual tensión).}
\comprobacion con $R_1=R_2$, la corriente se reparte a partes iguales.
\end{solucion}

\begin{solucion}{eje:5-minimos-cuadrados}
\datos puntos $(1,1)$, $(2,3)$, $(3,2)$.
\planteamiento resolver $D_a=D_b=0$ y comprobar con el \cref{teo:hessiano}.
\pasos
\begin{enumerate}[label=\arabic*.]
  \item $D_a=-2\sum x_i(y_i-ax_i-b)=0\iff a\sum x_i^2+b\sum x_i=\sum x_iy_i\iff14a+6b=13$.
  \item $D_b=-2\sum(y_i-ax_i-b)=0\iff a\sum x_i+3b=\sum y_i\iff6a+3b=6$.
  \item De la segunda, $b=2-2a$; sustituyendo, $14a+12-12a=13$, $a=\frac12$, $b=1$.
  \item $D_{aa}=2\sum x_i^2=28$, $D_{ab}=2\sum x_i=12$, $D_{bb}=6$; $\Delta=168-144=24>0$ y $D_{aa}>0$: mínimo.
\end{enumerate}
\resultado \textbf{$y=\frac12x+1$.}
\comprobacion los residuos son $1-1{,}5=-0{,}5$, $3-2=1$ y $2-2{,}5=-0{,}5$; suman~0, como exige $D_b=0$, y $D=1{,}5$.
\end{solucion}

\begin{solucion}{eje:5-degenerado}
\datos $f_1=x^4+y^4$, $f_2=-f_1$, $f_3=x^4-y^4$.
\planteamiento calcular gradiente y hessiana en el origen y estudiar el signo de $f_i$ directamente.
\pasos
\begin{enumerate}[label=\arabic*.]
  \item Las tres tienen gradiente $(\pm4x^3,\pm4y^3)$, nulo en el origen, y hessiana $\operatorname{diag}(\pm12x^2,\pm12y^2)$, nula en el origen: $\Delta=0$.
  \item $f_1\ge0=f_1(0,0)$: mínimo. $f_2\le0=f_2(0,0)$: máximo.
  \item $f_3(x,0)=x^4>0$ y $f_3(0,y)=-y^4<0$ para $x,y\neq0$: silla.
\end{enumerate}
\resultado \textbf{mínimo, máximo y silla, respectivamente, todos con $\Delta=0$: si $\Delta=0$ el criterio no permite concluir.}
\comprobacion en los tres casos la forma cuadrática $Q$ de~\eqref{eq:t5-taylor} es idénticamente nula, y el comportamiento lo deciden los términos de cuarto orden.
\end{solucion}

\begin{solucion}{eje:5-mle}
\datos $n_i>0$, $N=\sum n_i$; $\ell(\mathbf p)=\sum n_i\ln p_i$ con $\sum p_i=1$, $p_i>0$.
\planteamiento (a) Lagrange en el abierto $\{p_i>0\}$; (b) acotar $\ell(\mathbf p)-\ell(\hat{\mathbf p})$ con $\ln t\le t-1$.
\pasos
\begin{enumerate}[label=\arabic*.]
  \item (a) $\partial\ell/\partial p_i=n_i/p_i$ y $\nabla g=(1,\dots,1)$: $n_i/p_i=\lambda$, luego $p_i=n_i/\lambda$.
  \item $\sum p_i=1$ da $N/\lambda=1$, $\lambda=N$, y el único candidato es $\hat p_i=n_i/N$.
  \item (b) Para cualquier $\mathbf p$ admisible, $\ell(\mathbf p)-\ell(\hat{\mathbf p})=\sum n_i\ln\frac{p_i}{\hat p_i}\le\sum n_i\Bigl(\frac{p_i}{\hat p_i}-1\Bigr)$.
  \item Como $n_i/\hat p_i=N$, la cota es $\sum Np_i-\sum n_i=N-N=0$.
  \item La igualdad en $\ln t\le t-1$ solo se da para $t=1$, luego $\ell(\mathbf p)=\ell(\hat{\mathbf p})$ exige $\mathbf p=\hat{\mathbf p}$.
\end{enumerate}
\resultado \textbf{el máximo absoluto, único, es $\hat p_i=n_i/N$: la frecuencia relativa.}
\comprobacion con $\mathbf n=(2,1,1)$: $\hat{\mathbf p}=(\frac12,\frac14,\frac14)$ da $\ell=2\ln\frac12+2\ln\frac14\approx-4{,}159$, mientras que la uniforme da $4\ln\frac13\approx-4{,}394$, menor.
\end{solucion}

\begin{solucion}{eje:5-media-geometrica}
\datos $f=xyz$ en $T=\{x,y,z\ge0,\ x+y+z=S\}$, $S>0$.
\planteamiento Weierstrass para la existencia, Lagrange en el interior y especialización a $S=a+b+c$.
\pasos
\begin{enumerate}[label=\arabic*.]
  \item $T$ es cerrado y acotado y $f$ continua: existe el máximo. En el borde de $T$, $f=0$; en $(\frac S3,\frac S3,\frac S3)$, $f>0$. Luego el máximo está en el interior ($x,y,z>0$).
  \item Lagrange con $g=x+y+z$, $\nabla g=(1,1,1)\neq\vcero$: $yz=xz=xy=\lambda$.
  \item Con $z>0$, $yz=xz$ da $x=y$; con $x>0$, $xz=xy$ da $y=z$. Así $x=y=z=\frac S3$ y $f=(\frac S3)^3$.
  \item (b) Si $a+b+c=0$, los tres son nulos y la desigualdad es $0\le0$. Si no, $(a,b,c)\in T$ con $S=a+b+c$, luego $abc\le\bigl(\frac{a+b+c}{3}\bigr)^3$, y se toma la raíz cúbica (creciente).
\end{enumerate}
\resultado \textbf{$\max xyz=(S/3)^3$ en $x=y=z=S/3$; $\sqrt[3]{abc}\le\frac{a+b+c}{3}$.}
\comprobacion $(a,b,c)=(1,2,3)$: $\sqrt[3]6\approx1{,}817\le2$.
\end{solucion}

\begin{solucion}{eje:5-cs-lagrange}
\datos $f(\vx)=\va\cdot\vx$, $g(\vx)=\norm\vx^2=1$, $\va\neq\vcero$.
\planteamiento Lagrange sobre la esfera, que es cerrada y acotada.
\pasos
\begin{enumerate}[label=\arabic*.]
  \item $\nabla f=\va$, $\nabla g=2\vx$ (no nulo en la esfera): $\va=2\lambda\vx$. Como $\va\neq\vcero$, $\lambda\neq0$ y $\vx=\va/(2\lambda)$.
  \item $\norm\vx=1$ da $\abs{2\lambda}=\norm\va$: $\vx=\pm\va/\norm\va$, con $f=\pm\norm\va$.
  \item Por Weierstrass, el máximo es $\norm\va$ y el mínimo $-\norm\va$.
  \item (b) Para $\mathbf y\neq\vcero$, $\vx=\mathbf y/\norm{\mathbf y}$ está en la esfera: $\abs{\va\cdot\mathbf y}=\norm{\mathbf y}\abs{\va\cdot\vx}\le\norm{\mathbf y}\norm\va$. Si $\mathbf y=\vcero$ es trivial.
  \item (c) $f$ es lineal con gradiente constante $\va$, y $f(\vu)=D_{\vu}f$: es el \cref{cor:maximo-crecimiento}, cuyo máximo $\norm{\nabla f}$ se alcanza en la dirección del gradiente.
\end{enumerate}
\resultado \textbf{$\max=\norm\va$ en $\va/\norm\va$; $\min=-\norm\va$ en $-\va/\norm\va$; de ahí Cauchy--Schwarz.}
\comprobacion $\va=(1,2,2)$: en $\vx=\frac13(1,2,2)$, $f=\frac{9}{3}=3=\norm\va$.
\end{solucion}

% ---------------------------------------------------------------------
\section{Tema 6}

\begin{solucion}{eje:6-simetria}
\datos el disco unidad $D$, simétrico respecto de ambos ejes.
\planteamiento si el integrando es impar en $x$ (o en $y$), la integral es nula; si es par, se reduce a la mitad del disco.
\pasos
\begin{enumerate}[label=\arabic*.]
  \item $x$ es impar en $x$: $\iint_Dx\dA=0$. $x^2y$ y $xy$ son impares en $y$: integrales nulas.
  \item $e^x>0$: integral positiva, no nula. Es par en $y$: $\iint_De^x\dA=2\iint_{D,\,y\ge0}e^x\dA$.
  \item $x^2$ es par en $x$ y en $y$: $\iint_Dx^2\dA=4\iint_{D,\,x,y\ge0}x^2\dA$.
  \item $\iint_D(x^2+y)\dA=\iint_Dx^2\dA+0=4\iint_{D,\,x,y\ge0}x^2\dA$.
\end{enumerate}
\resultado \textbf{nulas: $x$, $x^2y$, $xy$. No nulas: $e^x$ (dos veces el semidisco superior), $x^2$ y $x^2+y$ (cuatro veces el cuarto de disco con integrando $x^2$).}
\comprobacion en polares, $\iint_Dx^2\dA=\int_0^{2\pi}\cos^2\theta\,d\theta\int_0^1r^3dr=\frac\pi4$, y el cuarto de disco da $\frac\pi{16}$, la cuarta parte.
\end{solucion}

\begin{solucion}{eje:6-iteradas}
\datos dos integrales iteradas.
\planteamiento integrar primero la variable interior.
\pasos
\begin{enumerate}[label=\arabic*.]
  \item (a) $\int_{-1}^1(6x^2+2y)\,dy=\bigl[6x^2y+y^2\bigr]_{-1}^{1}=12x^2$; $\int_0^212x^2dx=32$.
  \item (b) $\int_0^u\sqrt{u^2+4}\,dv=u\sqrt{u^2+4}$; $\int_0^1u\sqrt{u^2+4}\,du=\bigl[\frac13(u^2+4)^{3/2}\bigr]_0^1=\frac{5\sqrt5-8}{3}$.
\end{enumerate}
\resultado \textbf{(a) $32$; (b) $\frac{5\sqrt5-8}{3}\approx1{,}060$.}
\comprobacion (a) en el otro orden: $\int_{-1}^1\int_0^2(6x^2+2y)\,dx\,dy=\int_{-1}^1(16+4y)\,dy=32$.
\end{solucion}

\begin{solucion}{eje:6-rectangulos}
\datos tres integrandos sobre rectángulos.
\planteamiento \cref{teo:fubini}, eligiendo el orden.
\pasos
\begin{enumerate}[label=\arabic*.]
  \item (a) $\int_0^1(x+y)^2dx=\frac{(1+y)^3-y^3}{3}=\frac{1+3y+3y^2}{3}$; $\int_0^2\frac{1+3y+3y^2}3dy=\frac{2+6+8}{3}=\frac{16}3$.
  \item (b) $\int_0^{\pi/2}\sin(x+y)\,dx=\cos y-\cos(y+\frac\pi2)=\cos y+\sin y$; $\int_0^{\pi/2}(\cos y+\sin y)\,dy=1+1=2$.
  \item (c) Integrando primero en $y$: $\int_0^1xe^{xy}dy=\bigl[e^{xy}\bigr]_{y=0}^{1}=e^x-1$; $\int_0^1(e^x-1)\,dx=e-2$. (En el otro orden habría que integrar por partes.)
\end{enumerate}
\resultado \textbf{(a) $\frac{16}3$; (b) $2$; (c) $e-2\approx0{,}7183$.}
\comprobacion (a) desarrollando, $\iint(x^2+2xy+y^2)=\frac23+2+\frac83=\frac{16}{3}$, con $\iint x^2=\frac13\cdot2$, $\iint2xy=2\cdot\frac12\cdot2$ e $\iint y^2=1\cdot\frac83$.
\end{solucion}

\begin{solucion}{eje:6-planteo}
\datos (a) triángulo $(0,0)$, $(0,2)$, $(-2,2)$; (b) región bajo $y=2x-x^2$ sobre el eje $x$.
\planteamiento dibujar la región y leer los límites (\cref{def:regiones}).
\pasos
\begin{enumerate}[label=\arabic*.]
  \item (a) Los lados son $x=0$, $y=2$ y la recta $y=-x$ que une $(0,0)$ y $(-2,2)$. La región es $\{0\le y\le2,\ -y\le x\le0\}$.
  \item $dx\,dy$: $\int_0^2\int_{-y}^{0}f\,dx\,dy$. $dy\,dx$: $\int_{-2}^{0}\int_{-x}^{2}f\,dy\,dx$.
  \item (b) $2x-x^2\ge0$ en $[0,2]$, con máximo $1$ en $x=1$. $dy\,dx$: $\int_0^2\int_0^{2x-x^2}f\,dy\,dx$.
  \item Despejando $x$ de $y=2x-x^2$: $x=1\pm\sqrt{1-y}$. $dx\,dy$: $\int_0^1\int_{1-\sqrt{1-y}}^{1+\sqrt{1-y}}f\,dx\,dy$.
  \item Área (b): $\int_0^2(2x-x^2)\,dx=4-\frac83=\frac43$.
\end{enumerate}
\resultado \textbf{los límites anteriores; área de (b) $=\frac43$.}
\comprobacion en el otro orden, $\int_0^12\sqrt{1-y}\,dy=2\cdot\frac23=\frac43$; y el área de (a) sale $2$ en ambos órdenes.
\end{solucion}

\begin{solucion}{eje:6-region}
\datos $D$: triángulo de vértices $(0,0)$, $(2,0)$, $(0,1)$.
\planteamiento región de tipo~I: $0\le x\le2$, $0\le y\le(2-x)/2$.
\pasos
\begin{enumerate}[label=\arabic*.]
  \item $\int_0^{(2-x)/2}x\,dy=\frac{x(2-x)}{2}$.
  \item $\int_0^2\frac{2x-x^2}{2}dx=\frac12\bigl(4-\frac83\bigr)=\frac23$.
\end{enumerate}
\resultado \textbf{$\frac23$.}
\comprobacion área $=1$ y la abscisa del centroide de un triángulo es la media de las de los vértices, $\frac{0+2+0}{3}=\frac23$; $\iint x\dA=\bar x\cdot\text{área}=\frac23$.
\end{solucion}

\begin{solucion}{eje:6-volumen-xy}
\datos $z=xy$ sobre $\{0\le x\le1,\ x^2\le y\le x\}$.
\planteamiento volumen $=\iint xy\dA$ (el integrando es no negativo).
\pasos
\begin{enumerate}[label=\arabic*.]
  \item $\int_{x^2}^{x}xy\,dy=x\frac{x^2-x^4}{2}$.
  \item $\frac12\int_0^1(x^3-x^5)\,dx=\frac12\bigl(\frac14-\frac16\bigr)=\frac1{24}$.
\end{enumerate}
\resultado \textbf{$\frac1{24}$.}
\comprobacion en el orden $dx\,dy$, con $y\le x\le\sqrt y$: $\int_0^1y\frac{y-y^2}{2}dy=\frac12(\frac13-\frac14)=\frac1{24}$.
\end{solucion}

\begin{solucion}{eje:6-orden-1}
\datos $\int_0^2\int_x^2e^{-y^2}dy\,dx$.
\planteamiento la región $\{0\le x\le y\le2\}$ es de tipo~II: $0\le y\le2$, $0\le x\le y$.
\pasos
\begin{enumerate}[label=\arabic*.]
  \item $\int_0^2\int_0^ye^{-y^2}dx\,dy=\int_0^2ye^{-y^2}dy$.
  \item $=\bigl[-\frac12e^{-y^2}\bigr]_0^2=\frac{1-e^{-4}}{2}$.
\end{enumerate}
\resultado \textbf{$\frac{1-e^{-4}}{2}\approx0{,}4908$.}
\comprobacion el integrando es menor que $1$ y el triángulo tiene área $2$, luego la integral debe ser menor que $2$; y es positiva.
\end{solucion}

\begin{solucion}{eje:6-orden-2}
\datos $\int_0^1\int_{\sqrt[3]x}^{1}\frac{du\,dx}{1+u^4}$.
\planteamiento $\{0\le x\le1,\ x^{1/3}\le u\le1\}=\{0\le u\le1,\ 0\le x\le u^3\}$.
\pasos
\begin{enumerate}[label=\arabic*.]
  \item $\int_0^1\int_0^{u^3}\frac{dx}{1+u^4}\,du=\int_0^1\frac{u^3}{1+u^4}du$.
  \item $=\bigl[\frac14\ln(1+u^4)\bigr]_0^1=\frac{\ln2}4$.
\end{enumerate}
\resultado \textbf{$\frac{\ln2}{4}\approx0{,}1733$.}
\comprobacion la región tiene área $\int_0^1u^3du=\frac14$ y el integrando está entre $\frac12$ y $1$: $0{,}125\le0{,}1733\le0{,}25$.
\end{solucion}

\begin{solucion}{eje:6-polares}
\datos cuarto de disco de radio $a$; integrando $1/(1+r^2)$.
\planteamiento \cref{teo:polares} con $0\le\theta\le\frac\pi2$, $0\le r\le a$.
\pasos
\begin{enumerate}[label=\arabic*.]
  \item $\int_0^{\pi/2}\int_0^a\frac{r}{1+r^2}dr\,d\theta=\frac\pi2\cdot\bigl[\frac12\ln(1+r^2)\bigr]_0^a$.
  \item $=\frac\pi4\ln(1+a^2)$.
\end{enumerate}
\resultado \textbf{$\frac\pi4\ln(1+a^2)$.}
\comprobacion para $a$ pequeño, $\ln(1+a^2)\approx a^2$ y el resultado es $\approx\frac{\pi a^2}{4}$, el área del cuarto de disco, como corresponde a un integrando próximo a~1.
\end{solucion}

\begin{solucion}{eje:6-semiesfera}
\datos $z=\sqrt{a^2-r^2}$ sobre $r\le a$.
\planteamiento \cref{teo:polares}.
\pasos
\begin{enumerate}[label=\arabic*.]
  \item $V=\int_0^{2\pi}\int_0^a\sqrt{a^2-r^2}\,r\,dr\,d\theta$.
  \item Con $s=a^2-r^2$: $\int_0^a\sqrt{a^2-r^2}\,r\,dr=\bigl[-\frac13(a^2-r^2)^{3/2}\bigr]_0^a=\frac{a^3}{3}$.
  \item $V=2\pi\cdot\frac{a^3}3=\frac23\pi a^3$.
\end{enumerate}
\resultado \textbf{$\frac23\pi a^3$.}
\comprobacion es la mitad del volumen de la esfera, $\frac43\pi a^3$.
\end{solucion}

\begin{solucion}{eje:6-centroide}
\datos sector $\{0\le r\le a,\ -\alpha\le\theta\le\alpha\}$.
\planteamiento por simetría $\bar y=0$; calcular área y $\iint x\dA$ en polares.
\pasos
\begin{enumerate}[label=\arabic*.]
  \item Área $=\int_{-\alpha}^{\alpha}\int_0^ar\,dr\,d\theta=\alpha a^2$.
  \item $\iint x\dA=\int_{-\alpha}^{\alpha}\cos\theta\,d\theta\int_0^ar^2dr=2\sin\alpha\cdot\frac{a^3}3$.
  \item $\bar x=\frac{2a^3\sin\alpha/3}{\alpha a^2}=\frac{2a\sin\alpha}{3\alpha}$.
  \item Semidisco, $\alpha=\frac\pi2$: $\bar x=\frac{4a}{3\pi}\approx0{,}4244\,a$.
\end{enumerate}
\resultado \textbf{$(\bar x,\bar y)=\bigl(\frac{2a\sin\alpha}{3\alpha},0\bigr)$; semidisco: $\frac{4a}{3\pi}$.}
\comprobacion cuando $\alpha\to0$, $\bar x\to\frac{2a}{3}$: un sector muy estrecho es casi un triángulo, cuyo centroide está a $\frac23$ de la altura desde el vértice. Con $\alpha=\pi$ (disco completo), $\bar x=0$.
\end{solucion}

\begin{solucion}{eje:6-gauss-polar}
\datos $e^{-(x^2+y^2)}$ sobre el disco de radio $R$.
\planteamiento \cref{teo:polares}; el valor sobre todo el plano es $\pi$ (\cref{teo:gauss}).
\pasos
\begin{enumerate}[label=\arabic*.]
  \item $\int_0^{2\pi}\int_0^Re^{-r^2}r\,dr\,d\theta=2\pi\bigl[-\frac12e^{-r^2}\bigr]_0^R=\pi(1-e^{-R^2})$.
  \item $\pi(1-e^{-R^2})=0{,}99\pi\iff e^{-R^2}=0{,}01\iff R^2=\ln100$.
  \item $R=\sqrt{\ln100}\approx2{,}146$.
\end{enumerate}
\resultado \textbf{(a) $\pi(1-e^{-R^2})$; (b) $R=\sqrt{\ln100}\approx2{,}146$.}
\comprobacion $e^{-2{,}146^2}=e^{-4{,}605}\approx0{,}0100$.
\end{solucion}

\begin{solucion}{eje:6-capas}
\datos $z=f(r)\ge0$ continua en $[0,a]$; en (b), $f(r)=h(1-r/a)$.
\planteamiento la integral en $\theta$ es inmediata porque el integrando no depende de $\theta$.
\pasos
\begin{enumerate}[label=\arabic*.]
  \item (a) $V=\int_0^{2\pi}\int_0^af(r)\,r\,dr\,d\theta=\int_0^{2\pi}d\theta\cdot\int_0^arf(r)\,dr=2\pi\int_0^arf(r)\,dr$, que es la fórmula de las capas.
  \item La interpretación coincide: $2\pi r\,f(r)\,dr$ es el volumen de una capa cilíndrica de radio $r$, altura $f(r)$ y grosor $dr$, y el anillo polar de radio $r$ tiene área $2\pi r\,dr$.
  \item (b) $V=2\pi h\int_0^a\bigl(r-\frac{r^2}{a}\bigr)dr=2\pi h\bigl(\frac{a^2}2-\frac{a^2}3\bigr)=\frac13\pi a^2h$.
\end{enumerate}
\resultado \textbf{ambas fórmulas coinciden; el cono tiene volumen $\frac13\pi a^2h$.}
\comprobacion para el cilindro, $f\equiv h$: $2\pi h\frac{a^2}{2}=\pi a^2h$.
\end{solucion}

\begin{solucion}{eje:6-cota}
\datos $f(x,y)=\frac{x}{1+x^2+y^2}$ en $[0,1]^2$.
\planteamiento acotar por una función que solo dependa de $x$ y usar la monotonía (\cref{prop:propiedades-integral}(ii)) y el \cref{cor:separable}.
\pasos
\begin{enumerate}[label=\arabic*.]
  \item Para $x,y\in[0,1]$: $1+x^2+y^2\ge1+x^2>0$ y $x\ge0$, luego $f(x,y)\le\frac{x}{1+x^2}$.
  \item Por monotonía, $\iint f\dA\le\iint\frac{x}{1+x^2}\dA=\int_0^1\frac{x}{1+x^2}dx\cdot\int_0^1dy$.
  \item $\int_0^1\frac{x}{1+x^2}dx=\bigl[\frac12\ln(1+x^2)\bigr]_0^1=\frac{\ln2}{2}\approx0{,}34657$.
  \item Como $0{,}34657<0{,}35$, la integral es menor que $0{,}35$.
\end{enumerate}
\resultado \textbf{$\iint f\dA\le\frac{\ln2}2\approx0{,}3466<0{,}35$.}
\comprobacion el valor numérico de la integral es $\approx0{,}2878$, efectivamente por debajo de la cota.
\end{solucion}

\begin{solucion}{eje:6-gaussiana-normal}
\datos $\mu\in\R$, $\sigma>0$; \cref{teo:gauss,cor:gauss-normal}.
\planteamiento (a) cambio de variable; (b) partes en $[-R,R]$ y paso al límite; (c) imparidad y (b).
\pasos
\begin{enumerate}[label=\arabic*.]
  \item (a) Con $u=\frac{x-\mu}{\sigma\sqrt2}$, $dx=\sigma\sqrt2\,du$, y los límites $\pm\infty$ se conservan: $\int e^{-(x-\mu)^2/(2\sigma^2)}dx=\sigma\sqrt2\int e^{-u^2}du=\sigma\sqrt{2\pi}$.
  \item (b) En $[-R,R]$, con $u=x$ y $dv=xe^{-x^2/2}dx$ ($v=-e^{-x^2/2}$): $\int_{-R}^{R}x^2e^{-x^2/2}dx=\bigl[-xe^{-x^2/2}\bigr]_{-R}^{R}+\int_{-R}^{R}e^{-x^2/2}dx=-2Re^{-R^2/2}+\int_{-R}^{R}e^{-x^2/2}dx$.
  \item Cuando $R\to\infty$, $Re^{-R^2/2}\to0$ y la última integral tiende a $\sqrt{2\pi}$.
  \item (c) $\E[X]=\int x\varphi(x)\,dx$ converge (una primitiva de $\abs x e^{-x^2/2}$ en $x>0$ es $-e^{-x^2/2}$) y es nula porque el integrando es impar. $\E[X^2]=\frac{1}{\sqrt{2\pi}}\sqrt{2\pi}=1$.
\end{enumerate}
\resultado \textbf{(a) $\sigma\sqrt{2\pi}$; (b) $\sqrt{2\pi}$; (c) $\E[X]=0$, $\E[X^2]=1$, de modo que $\operatorname{Var}(X)=1$.}
\comprobacion con $\sigma=1$, $\mu=0$, (a) se reduce al \cref{cor:gauss-normal}. Numéricamente, $\int_{-5}^{5}x^2e^{-x^2/2}dx\approx2{,}5066\approx\sqrt{2\pi}$.
\end{solucion}
